\section{请求首先需要一个稳定的边界契约}

\subsection{只有入口地址还不能表达请求}

假设处理器已经提供一条指令，能把执行权从用户程序转交给内核的固定入口。此时仍然缺少
两个事实：用户究竟请求哪项服务，以及服务所需的参数在哪里。让每项服务占用一个不同入口
地址看似直接，却会让入口数量随着服务增加而增长；内核还要为每个入口单独配置权限和返回
规则。更重要的是，程序无法只凭一个稳定编号描述请求，兼容层、审计器和过滤器也难以统一
观察所有入口。

因此，最小设计只保留一个受控入口，并给每项服务分配一个整数编号。用户在进入前把编号和
参数放进约定位置；内核先读取编号，再选择处理函数。这个“编号放在哪里、参数怎样排列、
结果怎样返回”的约定属于\term{应用二进制接口}{application binary interface, ABI}。
这里先引入 ABI，是因为跨边界的两段代码可能由不同编译器、不同语言甚至不同年代生成，
它们不能依赖 C 函数名、类型系统或源代码声明来彼此理解。

一个教学用的最小请求可以表示为：

\begin{lstlisting}
boundary request:
  operation_number
  argument_0 ... argument_5
  user_continuation
  result_slot
\end{lstlisting}

\code{operation_number} 回答“请求什么”；参数回答“对哪个对象、使用哪些数值”；用户继续
位置回答“处理完回到哪里”；结果槽回答“内核如何把成功范围或失败原因交还给调用者”。这
四类状态不会因为采用寄存器或内存传递而消失。真实 ABI 只是为它们选择了具体载体。

\begin{longtable}{|L{0.20\linewidth}|L{0.28\linewidth}|L{0.40\linewidth}|}
\hline
边界信息 & 最小作用 & 若缺失会发生什么 \\ \hline
操作编号 & 从统一入口选择服务 & 内核不知道调用者希望修改哪类状态 \\ \hline
参数集合 & 描述对象、地址、长度和选项 & 处理函数只能执行固定动作，无法服务不同对象 \\ \hline
继续位置 & 保存用户程序下一步 & 内核完成后无法恢复原控制流 \\ \hline
结果或错误 & 说明已经完成的范围 & 调用者无法区分成功、短完成和失败 \\ \hline
\end{longtable}

\subsection{x86-64 为什么选用这些寄存器}

以常见的 x86-64 Linux ABI 为具体实例。进入前，\code{RAX} 保存系统调用号，六个整数或
指针参数依次放在 \code{RDI}、\code{RSI}、\code{RDX}、\code{R10}、\code{R8} 和
\code{R9}。返回值仍放在 \code{RAX}。第四个参数使用 \code{R10} 而不是普通函数 ABI
常用的 \code{RCX}，原因不是任意偏好：\code{SYSCALL} 指令会把用户返回地址写入
\code{RCX}，原有 \code{RCX} 内容不能继续承担第四参数。

\begin{longtable}{|L{0.17\linewidth}|L{0.27\linewidth}|L{0.44\linewidth}|}
\hline
寄存器 & 边界含义 & 进入后为什么仍要保存 \\ \hline
\code{RAX} & 系统调用号，返回时改为结果 & 原编号要保留在 \code{orig\_ax} 一类字段中，供审计、重启和跟踪使用 \\ \hline
\code{RDI--R9} & 最多六个直接参数 & 处理函数、过滤器或跟踪器可能在不同阶段读取它们 \\ \hline
\code{RCX} & \code{SYSCALL} 自动写入用户返回 \code{RIP} & 返回前要恢复用户下一条指令地址 \\ \hline
\code{R11} & \code{SYSCALL} 自动写入用户 \code{RFLAGS} & 返回时要恢复允许返回的用户标志 \\ \hline
\code{RSP} & 进入前仍指向用户栈 & \code{SYSCALL} 不自动保存或切换它，入口代码必须先保护其值 \\ \hline
\end{longtable}

这是一项体系结构相关的实例，不是所有系统的唯一布局。AArch64 常用 \code{x8} 保存系统
调用号、\code{x0--x5} 保存参数，RISC-V 常用 \code{a7} 保存编号、\code{a0--a5} 保存
参数。稳定的操作系统语义是“编号、参数、继续状态和结果必须有约定位置”；具体寄存器由
体系结构 ABI 决定。

\begin{pdfdiagram}
  \node[os state, minimum width=45mm, minimum height=48mm] (u) at (0,0) {};
  \node[os title] at (0,18mm) {用户执行现场};
  \node[os label, fill=none, text=BookInk] at (0,2mm)
    {\code{RAX = nr}\\\code{RDI..R9 = args}\\\code{RSP = user stack}\\\code{RIP = syscall}};

  \node[os control, minimum width=35mm, minimum height=28mm] (gate) at (60mm,0)
    {受控入口指令\\固定入口地址\\固定特权规则};

  \node[os memory, minimum width=48mm, minimum height=48mm] (k) at (124mm,0) {};
  \node[os title] at (124mm,18mm) {入口帧需要的语义};
  \node[os label, fill=none, text=BookInk] at (124mm,2mm)
    {原请求编号\\六个参数\\用户 RIP/RSP/flags\\返回值槽};

  \draw[os strong bus] (u.east) -- node[os label, above]{执行 \code{SYSCALL}} (gate.west);
  \draw[os strong bus] (gate.east) -- node[os label, above]{保存为稳定快照} (k.west);
\end{pdfdiagram}
\diagramnote{边界契约先于具体处理函数。寄存器是用户提交请求时的临时载体；入口帧把这些值转成内核可以在检查、阻塞和调度期间持续引用的状态。}

\subsection{返回值为什么要同时表达数量和错误}

仅用“成功/失败”一位无法描述许多真实结果。请求读取 4096 字节时，文件末尾可能只剩 600
字节；向管道写入 4096 字节时，信号到达前可能已经提交 1024 字节。若内核只返回“失败”，
调用者会重写已经完成的前缀；若只返回“成功”，调用者又会误以为全部完成。因此，字节流
接口通常用非负数表达已完成数量，用负的错误编号表达“没有可优先报告的完成前缀”。

x86-64 Linux 的原始系统调用约定把 \code{-EACCES}、\code{-EFAULT} 等负值放入
\code{RAX}。用户态 C 库包装层通常把它转换为返回 \code{-1}，并把正的错误编号写入线程
局部的 \code{errno}。这两个边界要分开理解：内核 ABI 返回负错误编号，C 库 API 再提供
符合 C 语言习惯的表示。

\begin{longtable}{|L{0.18\linewidth}|L{0.26\linewidth}|L{0.46\linewidth}|}
\hline
原始结果 & 调用者应如何解释 & 状态含义 \\ \hline
\code{4096} & 全部推进 & 本次请求的 4096 字节已达到该接口的完成边界 \\ \hline
\code{600} & 短完成 & 前 600 字节已提交；剩余部分要由调用者决定是否继续 \\ \hline
\code{0} & 依接口解释 & \code{read} 可表示 EOF；其他接口可能表示成功但无数量 \\ \hline
\code{-EINTR} & 可按接口规则重试 & 信号在没有更优先的部分完成结果时打断等待 \\ \hline
\code{-EFAULT} & 修正用户地址 & 参数复制没有形成可报告的成功前缀 \\ \hline
\end{longtable}

结果槽因而也是提交账本。处理函数不能先修改对象偏移 4096 字节，再在后半段故障时只返回
\code{-EFAULT}，否则调用者不知道哪些字节已经生效。后文讨论用户内存时会继续沿这条
不变量分析部分复制。

\subsection{结构体参数为什么需要长度与版本}

六个寄存器足以传递少量整数，却不适合直接承载可扩展配置。常见做法是让一个参数指向用户
结构体。结构体一旦成为 ABI，字段偏移就可能被旧程序长期依赖；在中间插入字段会改变后续
偏移，扩大指针或时间类型也会改变布局。边界结构因此常带 \code{size}、\code{version}、
flags 和保留字段。

\begin{lstlisting}
struct request_v1 {
  uint32_t size;
  uint32_t flags;
  uint64_t object_id;
  uint64_t user_buffer;
  uint64_t length;
  uint64_t reserved[3];
};
\end{lstlisting}

\code{size} 让内核知道用户实际提供到哪个字段；保留字段要求用户清零，给未来扩展留下
可验证空间；未知 flags 必须拒绝，而不是静默忽略可能改变安全语义的请求。内核把已知前缀
复制到清零后的内部结构，再验证范围和组合。这样，新内核能识别旧结构，旧内核也能明确
拒绝自己不理解的新语义。

\begin{lstlisting}
copy_versioned_request(user_ptr):
  copy fixed header containing size and flags
  reject size smaller than mandatory prefix
  reject unsupported flag bits
  zero initialize kernel request
  copy min(user_size, known_size) bytes
  require unknown tail or reserved fields to be zero
  validate every length, offset and object reference
\end{lstlisting}

这里的内部结构不是用户结构的永久别名。复制完成后，内核拥有一份不随用户线程并发修改的
快照；内部字段还可以使用内核指针、引用对象和已检查长度，而不受 ABI 对齐规则限制。

\section{用户控制流怎样变成可恢复的入口现场}

\subsection{处理器自动完成的动作其实很少}

现在已有编号和参数，但入口仍不能安全调用普通内核函数。进入前的 \code{RSP} 指向用户
栈，\code{RCX} 和 \code{R11} 即将被覆盖，其他通用寄存器仍含用户值。若入口代码立即
\code{push}，它会把内核状态写到用户可控制的地址；该页可能未映射，也可能被另一线程
改写。系统必须先区分处理器自动建立的最小状态与入口软件随后建立的完整状态。

在典型 x86-64 配置下，\code{SYSCALL} 的关键动作是：

\begin{enumerate}
\tightlist
\item 把下一条用户指令地址写入 \code{RCX}。
\item 把用户 \code{RFLAGS} 写入 \code{R11}，并按 \code{IA32\_FMASK} 清除指定标志。
\item 从预先配置的 MSR 取得内核入口 \code{RIP} 和代码段相关属性。
\item 进入内核权限级并从入口 \code{RIP} 继续。
\end{enumerate}

它不会自动把全部寄存器压入某个帧，也不会像某些异常门那样自动把用户 \code{RSP} 换成
任务内核栈。用户 \code{RSP} 仍然只存在于寄存器里。入口最早的几条指令必须在不依赖用户
栈的条件下暂存它，再取得当前任务的内核栈顶。

\begin{longtable}{|L{0.20\linewidth}|L{0.30\linewidth}|L{0.38\linewidth}|}
\hline
状态 & 谁首先保存 & 为什么不能省略 \\ \hline
用户 \code{RIP} & 处理器写入 \code{RCX} & 决定系统调用返回后的下一条用户指令 \\ \hline
用户 \code{RFLAGS} & 处理器写入 \code{R11} & 保留用户可恢复标志，同时阻止危险标志影响入口 \\ \hline
用户 \code{RSP} & 入口软件写入每 CPU 暂存区或等价位置 & \code{SYSCALL} 不自动保存，换栈前又不能使用用户栈 \\ \hline
参数与通用寄存器 & 入口软件写入内核栈帧 & C 处理函数、信号与重启路径需要稳定快照 \\ \hline
内核继续位置 & 普通调用链与调度切换帧保存 & 任务在内核中阻塞后必须从原等待点继续 \\ \hline
\end{longtable}

\subsection{最早入口为何依赖每 CPU 状态}

换栈前不能从当前用户栈读取内核对象。入口却必须找到“当前 CPU 正在运行哪个 task”以及
“这个 task 的内核栈顶在哪里”。一种常见组织是在每个 CPU 的受保护区域保存当前 task
指针、入口临时栈指针和用户 \code{RSP} 暂存槽。x86-64 内核还可能使用
\code{SWAPGS} 在用户与内核的 GS 基址之间切换，使固定偏移能够访问这组每 CPU 数据。

\begin{lstlisting}
per_cpu_entry_state:
  current_task
  kernel_stack_top
  saved_user_rsp
  entry_scratch
  cpu_id
\end{lstlisting}

这些字段属于当前处理器的入口执行环境，不等于 task 自身。CPU A 的
\code{current_task} 会随着调度改变；task T 的内核栈则随 T 的生命周期存在，T 迁移到
CPU B 后仍使用自己的内核栈。把两者混为一体会产生一个常见误解：好像内核栈属于 CPU。
实际上，每 CPU 状态只负责在入口最早阶段找到当前 task，长期保存阻塞调用链的是线程
私有内核栈。

\begin{pdfdiagram}
  \node[os state, minimum width=42mm, minimum height=42mm] (cpu) at (0,10mm) {};
  \node[os title] at (0,25mm) {CPU 入口状态};
  \node[os label, fill=none, text=BookInk] at (0,9mm)
    {\code{current = T7}\\\code{kstack\_top = K7}\\临时 \code{user\_rsp}};

  \node[os memory, minimum width=46mm, minimum height=50mm] (task) at (62mm,10mm) {};
  \node[os title] at (62mm,29mm) {task T7};
  \node[os label, fill=none, text=BookInk] at (62mm,10mm)
    {调度状态\\地址空间引用\\凭据引用\\内核栈边界};

  \node[os control, minimum width=46mm, minimum height=50mm] (stack) at (126mm,10mm) {};
  \node[os title] at (126mm,29mm) {T7 的内核栈};
  \node[os label, fill=none, text=BookInk] at (126mm,10mm)
    {入口帧\\C 调用链\\等待位置\\切换帧};

  \draw[os strong bus] (cpu.east) -- node[os label, above]{定位当前 task} (task.west);
  \draw[os strong bus] (task.east) -- node[os label, above]{取得栈顶} (stack.west);
\end{pdfdiagram}
\diagramnote{每 CPU 入口状态只解决“此刻是谁在这颗 CPU 上进入内核”；task 和它的内核栈解决“这条执行流阻塞或迁移后怎样继续”。}

\subsection{为什么入口帧不能只保存返回地址}

若处理函数永不阻塞、没有信号、没有调试、没有审计，也不调用会改寄存器的函数，保存
\code{RCX} 和 \code{R11} 似乎足够。真实内核不具备这些前提。处理函数遵循内核函数
调用约定，会覆盖调用者保存寄存器；page fault 或中断可能嵌套；任务可能在等待数据时被
调度出去；信号处理需要修改用户返回现场；调试器甚至可能观察或修改参数与结果。

入口因此在内核栈上建立一个稳定的用户寄存器快照。Linux 常用
\code{struct pt\_regs} 一类结构承载它。下面是按语义分组的教学布局，不承诺某个内核
版本的精确字节偏移：

\begin{longtable}{|L{0.20\linewidth}|L{0.28\linewidth}|L{0.40\linewidth}|}
\hline
字段组 & 代表性字段 & 使用者与必要性 \\ \hline
请求身份 & \code{orig\_ax} & 保留原系统调用号，重启、seccomp、审计和 ptrace 需要它 \\ \hline
直接参数 & \code{di, si, dx, r10, r8, r9} & 分发包装层从稳定帧读取，而不是依赖稍后已变化的物理寄存器 \\ \hline
通用现场 & \code{r15...rbx, bp} & 信号、调试、异常嵌套和返回路径需要完整用户可见状态 \\ \hline
返回结果 & \code{ax} & 处理函数写入成功数量或负错误编号 \\ \hline
控制现场 & \code{ip, cs, flags, sp, ss} & 验证并恢复用户指令、栈和权限相关状态 \\ \hline
\end{longtable}

\code{orig\_ax} 与 \code{ax} 必须分开。进入时两者都可源于用户 \code{RAX}；分发后
\code{ax} 被改成结果，\code{orig\_ax} 仍保存原编号。若信号路径决定重启调用，它需要
知道原请求而不是把返回值误当成编号。

\subsection{入口帧和调度切换帧为何是两类结构}

入口帧回答“怎样回到用户态”；调度切换帧回答“怎样继续一条内核控制流”。任务 T 在
\code{read} 中睡眠时，入口帧位于 T 的内核栈较高位置，保存最初用户状态；更靠近当前
栈顶的位置保存内核函数调用链和调度器切换所需的被调用者保存寄存器。调度器切到任务 U
时，不需要把 T 的整个 \code{pt\_regs} 复制到 U；它只保存足以恢复 T 内核栈指针和
内核继续地址的切换现场。

\begin{pdfdiagram}
  \node[os memory, minimum width=64mm, minimum height=78mm] (stack) at (0,0) {};
  \node[os title] at (0,33mm) {task T 的内核栈（高地址在上）};
  \draw[os guide] (-27mm,20mm) -- (27mm,20mm);
  \draw[os guide] (-27mm,-5mm) -- (27mm,-5mm);
  \draw[os guide] (-27mm,-25mm) -- (27mm,-25mm);
  \node[os label, fill=none, text=BookInk] at (0,27mm) {入口帧 \code{pt\_regs}};
  \node[os label, fill=none, text=BookInk] at (0,8mm) {\code{sys\_read} 与 VFS 调用链};
  \node[os label, fill=none, text=BookInk] at (0,-15mm) {等待与 \code{schedule} 调用链};
  \node[os label, fill=none, text=BookInk] at (0,-31mm) {调度切换帧与当前 \code{RSP}};

  \node[os state, minimum width=58mm, minimum height=28mm] (user) at (96mm,22mm)
    {入口帧的恢复目标\\用户 \code{RIP/RSP/RFLAGS}\\参数与返回值};
  \node[os control, minimum width=58mm, minimum height=28mm] (kernel) at (96mm,-24mm)
    {切换帧的恢复目标\\内核 \code{RSP}\\等待函数之后的继续位置};

  \draw[os strong bus] (stack.east |- 0,27mm) -- (user.west);
  \draw[os strong bus] (stack.east |- 0,-31mm) -- (kernel.west);
\end{pdfdiagram}
\diagramnote{一次阻塞系统调用同时保留两种继续关系。调度先用切换帧恢复内核控制流；处理函数结束后，返回路径才使用入口帧恢复用户控制流。}

这一区分也解释了为什么“系统调用导致上下文切换”不是必然事实。进入内核只建立入口现场；
若处理函数立即完成，仍由同一 task 在同一 CPU 上返回，不发生任务调度切换。只有调用
阻塞、时间片耗尽或更高优先级任务需要运行时，调度器才另行保存内核切换现场并选择其他
task。

\subsection{从用户指令到完整入口帧的逐步状态}

把前述局部结构合在一起，可以得到入口的第一条完整但尚未分发的 trace。假设 task T7
请求 \code{getpid}，系统调用号已放入 \code{RAX}，用户下一条指令地址为
\code{0x40102a}，用户栈顶为 \code{0x7fff...e800}。

\begin{longtable}{|L{0.12\linewidth}|L{0.22\linewidth}|L{0.28\linewidth}|L{0.28\linewidth}|}
\hline
阶段 & 当前栈 & 新形成的状态 & 尚未发生的动作 \\ \hline
执行入口前 & 用户栈 & \code{RAX=nr}，参数寄存器有效 & 未进入内核，未检查请求 \\ \hline
\code{SYSCALL} 后 & \code{RSP} 仍是用户值 & \code{RCX=0x40102a}，\code{R11=user flags} & 未自动保存其余寄存器，未自动换栈 \\ \hline
入口最早阶段 & 不使用用户栈 & 暂存用户 \code{RSP}，取得 T7 与 K7 栈顶 & 尚未调用 C 代码 \\ \hline
切到 K7 后 & T7 内核栈 & 按约定压入完整入口帧，写入 \code{orig\_ax} & 尚未根据编号选择服务 \\ \hline
入口帧完成 & T7 内核栈 & 内核已有稳定用户状态快照 & 尚未证明参数、对象或权限合法 \\ \hline
\end{longtable}

入口帧完成只是把不稳定寄存器转为稳定的内核记录，不代表请求已经获准。下一步才是从
\code{orig\_ax} 找到候选处理函数，并让过滤、参数复制和权限规则逐层缩小请求的可执行
范围。

\begin{keyidea}
受控入口首先是一项状态保存协议。处理器只提供最小权限转换和两个返回寄存器；入口软件必须
在不信任用户栈的前提下找到当前 task、切到线程内核栈，并构造能够经历函数调用、阻塞、
调度、信号和调试的入口帧。入口成功不等于请求获准，更不等于操作完成。
\end{keyidea}

\section{入口帧之后为何还要分发与逐层验证}

\subsection{服务增多以后需要从编号找到实现}

最初只有“读取进程编号”和“调整系统时钟”两项服务时，入口可以用两个条件分支选择实现。
当服务增长到数百项，连续条件分支既难维护，也无法直接验证编号是否落在合法范围。编号
已经是一个稠密或近似稠密的整数空间，因此可以把“编号到处理函数”的关系保存成数组。
这张数组称为系统调用分发表。

\begin{lstlisting}
syscall_table:
  [NR_read]          -> sys_read
  [NR_write]         -> sys_write
  [NR_getpid]        -> sys_getpid
  [NR_clock_settime] -> sys_clock_settime
\end{lstlisting}

分发表保存的是候选入口，不是授权清单。编号合法只说明内核认识这项操作；当前主体是否有权
执行、参数是否有效、目标对象是否处于允许状态，都要由后续路径判断。若把“表项存在”误当
成“请求获准”，任何进程都能调用修改时钟、挂载文件系统或控制网络配置的服务。

\begin{longtable}{|L{0.20\linewidth}|L{0.30\linewidth}|L{0.38\linewidth}|}
\hline
分发阶段检查 & 能证明什么 & 不能证明什么 \\ \hline
编号非负且小于表长 & 访问分发表不会越界 & 该编号在当前 ABI 中一定实现 \\ \hline
表项不是空缺处理函数 & 内核有候选实现 & 调用者具备对象权限 \\ \hline
ABI 类型匹配 & 参数能按正确位宽和顺序解释 & 参数值、用户地址和长度合法 \\ \hline
入口来源是用户态 & 当前帧可按用户返回规则处理 & 用户提交的内容可信 \\ \hline
\end{longtable}

越界或未实现编号通常返回 \code{-ENOSYS}。内核不能让用户编号直接变成未检查的函数指针
或数组越界索引。投机执行缓解还可能在范围检查后增加索引约束，防止处理器短暂沿错误路径
读取越界表项；这是在基本范围检查正确之后，为微体系结构信息泄露增加的防护。

\subsection{为什么分发包装层要从入口帧取参数}

入口汇编已经把参数保存进 \code{pt\_regs}。分发层可以按 ABI 从帧中取出六个值，再调用
类型化包装函数。这样，体系结构相关的寄存器排列集中在边界层，主体实现只面对普通整数、
句柄和用户地址标记。

\begin{lstlisting}
dispatch(regs):
  nr = regs.orig_ax
  if nr outside implemented range:
    regs.ax = -ENOSYS
    return
  args = decode_abi_arguments(regs)
  regs.ax = syscall_table[nr](args)
\end{lstlisting}

这里的 \code{args} 仍然只是未经信任的数值。若某个值表示长度，处理函数要检查上限和
加法溢出；若表示 fd，要取得稳定文件引用；若表示用户指针，要复制或固定页面；若表示
flags，要拒绝未知位。包装层完成“按 ABI 解码”，具体实现完成“按对象语义验证”，两者
不能由一次统一检查替代。

\subsection{入口过滤为什么位于实际处理函数之前}

有些进程只应使用系统调用集合中的一小部分。例如，负责解析不可信图片的沙箱进程可能只需
读写已有 fd、分配内存和退出，不需要创建网络连接或加载内核模块。若等到处理函数已经取得
对象锁、分配资源之后再拒绝，撤销成本高，也容易遗漏副作用。因此内核可以在分发前后设置
一组入口工作项：

\begin{enumerate}
\tightlist
\item 根据 task 的跟踪状态决定是否通知调试器。
\item 让 seccomp 过滤器观察系统调用号和原始参数，决定允许、返回错误、通知监管者或终止。
\item 建立审计上下文，记录主体、编号和稍后形成的结果。
\item 在检查通过后调用实际处理函数。
\end{enumerate}

\begin{pdfdiagram}[scale=0.88]
  \node[os state, minimum width=35mm] (frame) at (0,0)
    {入口帧\\编号与原始参数};
  \node[os control, minimum width=34mm] (filter) at (46mm,0)
    {入口策略\\seccomp/ptrace};
  \node[os compute, minimum width=34mm] (dispatch) at (92mm,0)
    {ABI 分发\\范围与表项};
  \node[os memory, minimum width=35mm] (impl) at (138mm,0)
    {对象实现\\复制与权限};

  \draw[os strong bus] (frame.east) -- (filter.west);
  \draw[os strong bus] (filter.east) -- (dispatch.west);
  \draw[os strong bus] (dispatch.east) -- (impl.west);
  \draw[os feedback] (filter.south) -- ++(0,-15mm) -|
    node[os label, below, pos=0.55]{拒绝：尚未产生对象副作用} (frame.south);
\end{pdfdiagram}
\diagramnote{入口过滤处理原始请求面，对象权限处理已解析的目标。前者适合收缩可调用操作集合，后者才能判断“这个主体是否能操作这个对象”。}

seccomp 看到的指针通常只是一个数值，不能安全遍历任意复杂对象语义；LSM 钩子则位于目标
对象已解析的位置，可以检查 inode、socket、task 等内核对象。把两者放在不同阶段不是
重复，而是因为它们拥有的稳定信息不同。

\subsection{最小 getpid 如何完成而不触碰用户内存}

\code{getpid} 适合作为第一条完整分发例子，因为它没有用户指针，也通常不会阻塞。进入
前，用户包装函数只提交系统调用号。入口建立帧后，分发层找到处理函数；处理函数从当前
task 的进程标识关系中取得对调用者可见的 PID；结果写入 \code{regs.ax}，随后进入返回
检查。

\begin{longtable}{|L{0.12\linewidth}|L{0.24\linewidth}|L{0.28\linewidth}|L{0.26\linewidth}|}
\hline
步骤 & 读取状态 & 写入状态 & 边界含义 \\ \hline
1 & 用户 \code{RAX=NR\_getpid} & \code{RCX/R11} 与入口临时状态 & 权限转换完成，请求尚未分发 \\ \hline
2 & 每 CPU \code{current} & T7 内核栈上的入口帧 & 用户现场成为稳定快照 \\ \hline
3 & \code{orig\_ax} & 选择 \code{sys\_getpid} & 编号已解析，但结果尚未形成 \\ \hline
4 & T7 的 PID namespace 关系 & \code{regs.ax=visible\_pid} & 服务结果形成 \\ \hline
5 & 返回工作标志与入口帧 & 用户 \code{RAX=visible\_pid} & 用户重新获得控制权 \\ \hline
\end{longtable}

即便这条路径不复制用户内存，返回值仍可能相对 PID namespace 解释。处理函数不能简单
返回一个不考虑调用者视图的全局整数。这个例子提前暴露一项通用规则：请求结果经常依赖
当前 task 所在的命名空间或凭据视图，入口帧本身不会替处理函数完成这类对象语义。

\subsection{检查顺序为什么必须服从信息依赖}

系统调用常有多项检查，但顺序不能只按代码整洁程度排列。某项检查若需要稳定对象，就必须
在取得对象引用之后；某项检查若能在分配资源前拒绝，就应尽量提前；某项检查若会泄露对象
是否存在，还要遵守接口规定的错误优先级。

以“通过 fd 写数据”为例，可以先形成如下依赖，而不急于展开文件系统：

\begin{enumerate}
\tightlist
\item 验证长度能否在内部类型中表示，避免后续计算溢出。
\item 用 fd 槽位取得稳定的打开对象引用；失败时返回 \code{-EBADF}。
\item 检查该打开对象是否允许写以及对象当前状态是否接受写入。
\item 把用户数据转为内核可控制的字节或页面引用。
\item 让对象实现提交尽可能长的前缀，并返回实际数量。
\item 释放临时页面、缓冲区和对象引用。
\end{enumerate}

有些实现会为性能调整复制与对象检查的具体先后，但必须保持两个不变量：在对象使用期间
持有稳定引用；在报告错误时，已经提交的前缀不能从结果中消失。后续小节分别推导这两项
结构。

\section{用户地址为何必须转换成内核拥有的状态}

\subsection{一个地址数值同时缺少三类保证}

用户把 \code{0x7fff...1000} 放入参数寄存器时，内核得到的只是一个整数。这个整数没有
附带三项保证：

\begin{enumerate}
\tightlist
\item 当前地址空间是否存在覆盖它的合法区间。
\item 对应页面此刻是否已建立、是否允许所需方向的访问。
\item 从检查到使用期间，另一个线程是否会改写映射或内容。
\end{enumerate}

第一项属于地址范围与 VMA 策略，第二项属于当前页表和缺页处理，第三项属于并发快照与
对象生命周期。仅检查“地址小于用户空间上界”最多排除明显的内核地址，不能证明页面已经
可读，更不能固定内容。

\begin{longtable}{|L{0.20\linewidth}|L{0.30\linewidth}|L{0.38\linewidth}|}
\hline
动作 & 得到的保证 & 仍然缺少的保证 \\ \hline
\code{access\_ok(ptr,n)} & 范围形式上位于允许的用户地址域且不溢出 & 页面存在、权限成立、内容稳定 \\ \hline
逐页访问或复制 & 每个实际访问字节要么成功，要么进入可修复 fault 路径 & 复制后的用户原位置不再变化 \\ \hline
复制到内核缓冲区 & 内核拥有当时的内容快照 & 原用户内容随后是否改变；通常已不再重要 \\ \hline
固定用户页 & 页面在异步期间不会被回收或换成另一页 & 用户是否还能改页内容、设备访问何时结束 \\ \hline
\end{longtable}

\subsection{先检查再直接解引用为何仍然有竞态}

假设线程 A 先验证用户页可读，随后在解析结构时反复从用户地址取字段；线程 B 可以在两次
读取之间改写长度或指针。A 第一次看见长度 16，分配 16 字节；第二次又从用户地址读到
长度 4096，再复制 4096 字节，就会越过已分配缓冲区。这类“检查时与使用时不一致”称为
\term{检查—使用时间竞争}{time of check to time of use, TOCTOU}。

\begin{lstlisting}
unsafe:
  if user_header.length <= 16:
    kbuf = allocate(user_header.length)
    copy(kbuf, user_header.data, user_header.length)
\end{lstlisting}

修复的关键不是增加一次相同检查，而是改变状态所有权：先把固定大小的头复制到内核局部
变量，所有后续判断都读取这份快照；长度通过检查后，再复制对应数据。若协议允许用户在
操作期间并发更新共享内存，就要使用明确的序列号、原子字段或锁协议，而不是假设普通读取
构成快照。

\begin{lstlisting}
safer:
  khdr = zero_initialized_header()
  copy_fixed_header_from_user(&khdr, user_header)
  validate(khdr.length, khdr.flags)
  kbuf = allocate(khdr.length)
  copy_bytes_from_user(kbuf, khdr.data, khdr.length)
\end{lstlisting}

\subsection{跨页复制为什么可能在中间故障}

请求复制 5000 字节时，起始地址可能落在某页最后 1000 字节，剩余 4000 字节位于下一页。
第一页可读并不能证明第二页可读。第二页可能尚未建立 PTE，但 VMA 允许按需分配；此时
page fault 可以建立页面后重试。也可能根本没有覆盖第二页的 VMA；此时复制必须收束为
错误或部分完成。

\begin{pdfdiagram}[scale=0.88]
  \node[os memory, minimum width=44mm, minimum height=34mm] (p0) at (0,12mm)
    {用户页 P0\\末尾 1000 B\\可读、已映射};
  \node[os memory, minimum width=44mm, minimum height=34mm] (p1) at (0,-26mm)
    {用户页 P1\\前 4000 B\\PTE 尚不存在};

  \node[os compute, minimum width=46mm] (copy0) at (61mm,12mm)
    {复制第一段\\\code{copied=1000}};
  \node[os control, minimum width=46mm] (fault) at (61mm,-26mm)
    {访问第二段\\产生 page fault};

  \node[os state, minimum width=51mm, minimum height=48mm] (decide) at (126mm,-7mm) {};
  \node[os title] at (126mm,10mm) {缺页决策};
  \node[os label, fill=none, text=BookInk] at (126mm,-7mm)
    {VMA 允许：建立页后重试\\VMA 不允许：转 fixup\\报告未复制范围};

  \draw[os strong bus] (p0.east) -- (copy0.west);
  \draw[os strong bus] (p1.east) -- (fault.west);
  \draw[os strong bus] (copy0.east) -- (decide.west |- copy0.east);
  \draw[os strong bus] (fault.east) -- (decide.west |- fault.east);
\end{pdfdiagram}
\diagramnote{用户复制以实际访问为准逐页推进。范围检查不是访问保证；可修复缺页与不可访问地址必须在 fault 路径中分开。}

\subsection{异常表怎样把预期 fault 变成 EFAULT}

普通内核指针访问坏地址通常说明内核错误。用户复制却有意访问不可信地址，fault 是接口
必须处理的输入结果。内核需要一种方法区分这两类内核态 fault。常见方案是在异常表中记录
“复制指令地址 → 修复地址”的映射。

\begin{lstlisting}
copy instruction at A:
  load from user address

exception table entry:
  fault_ip A -> fixup_ip F

page fault handler:
  if fault cannot be resolved and current ip has fixup:
    set instruction pointer to F
    continue with "bytes not copied" result
  else:
    follow ordinary kernel fault policy
\end{lstlisting}

异常表不让错误页面 magically 变得可读。它只提供一个受控退出点，使复制 helper 能报告
剩余字节，而不是把预期的坏用户参数升级成内核崩溃。上层系统调用再根据自己的提交语义，
把 helper 结果转换为 \code{-EFAULT}、短完成数量或其他接口结果。

\subsection{部分复制和部分系统调用完成不是同一层结果}

许多底层复制 helper 返回“还有多少字节未复制”，而系统调用返回“已完成多少字节”或负
错误。两者方向不同，且中间还隔着对象提交。例如，内核成功从用户地址复制 1000 字节，
不代表文件对象已经接受这 1000 字节；相反，对某些分段写路径，前一批字节可能已经提交，
后一批复制才遇到 fault。

\begin{longtable}{|L{0.18\linewidth}|L{0.24\linewidth}|L{0.28\linewidth}|L{0.20\linewidth}|}
\hline
时刻 & 用户复制进度 & 对象提交进度 & 可返回结果 \\ \hline
尚未复制 & 0/5000 & 0/5000 & 若立即失败，返回负错误 \\ \hline
内核快照完成 & 5000/5000 & 0/5000 & 仍可能因对象状态失败 \\ \hline
第一段提交 & 5000/5000 & 1000/5000 & 后续失败时通常优先返回 1000 \\ \hline
全部提交 & 5000/5000 & 5000/5000 & 返回 5000 \\ \hline
\end{longtable}

设计处理函数时，应分别维护 \code{copied\_from\_user} 与
\code{committed\_to\_object}，不能用一个计数器同时代表两种进度。最终返回值通常由对象
提交进度决定，因为它告诉调用者重试时应从哪里继续。

\subsection{输出参数为什么要先在内核中完整构造}

有些系统调用把结构体写回用户地址。若内核逐字段直接写用户结构，写到一半遇到坏页，用户
可能观察到新旧字段混合；若结构体含未初始化 padding，还可能泄露内核栈内容。稳妥路径是
先在清零的内核对象中构造完整结果，设置明确版本和长度，再进行一次边界复制。

\begin{lstlisting}
produce_user_result(user_out):
  result = zero_initialized()
  fill documented fields
  result.size = documented_size
  validate output range
  if copy_to_user(user_out, &result, result.size) fails:
    return -EFAULT
  return 0
\end{lstlisting}

如果接口规定输出可以部分形成，就必须明确每个字段的有效条件或返回实际数量。默认情况下，
让用户看到一个完整旧结果或完整新结果，比暴露半个结构更容易形成稳定契约。

\subsection{嵌套指针为什么要逐层形成快照}

\code{execve} 的 \code{argv}、\code{sendmsg} 的 iovec、若干 \code{ioctl} 结构都含
用户指针数组。复制外层结构只冻结了指针数值，没有冻结每个指针指向的内容。内核要先限制
数组个数和总长度，复制指针数组，再逐项复制或固定实际数据；每次加法和乘法都要检查溢出。

\begin{lstlisting}
copy_iovecs(user_iov, count):
  reject count above interface limit
  allocate zeroed kernel iovec array
  copy outer array once
  total = 0
  for each copied entry:
    reject length or total + length overflow
    validate address range for requested direction
    total += length
  return stable kernel array and checked total
\end{lstlisting}

外层数组、内层数据和总长度是三个不同状态。异步接口若在系统调用返回后仍使用内层页面，
还必须把页面固定和解除固定的责任交给长期请求对象；单纯保存用户指针不延长页面生命周期。

\begin{keyidea}
用户指针不是内核对象引用。范围检查只排除明显非法数值，实际复制还要处理逐页 fault；
复制到内核缓冲区才能形成内容快照，固定页面才能延长页面身份，但仍不自动禁止用户修改。
系统调用返回值必须依据对象已提交范围，而不是把底层复制计数直接照搬给用户。
\end{keyidea}

\section{对象引用怎样封闭检查与使用之间的竞态}

\subsection{fd 整数为何不能作为长期对象身份}

用户传入 fd 3 时，内核不能把整数 3 保存到异步请求里，稍后再查一次。线程 B 可能在此
期间关闭 fd 3，另一次 \code{open} 又复用同一槽位。稍后查询得到的可能是完全不同的对象。
因此，系统调用在开始使用对象时要从 fd table 取得指向打开文件对象的稳定引用，并增加
引用计数；后续路径持有对象引用，而不是反复解释整数。

\begin{pdfdiagram}
  \node[os state, minimum width=38mm] (fd0) at (0,16mm)
    {时刻 A\\fd 3};
  \node[os memory, minimum width=43mm] (file0) at (54mm,16mm)
    {打开对象 F\\\code{refcount=2}};
  \node[os control, minimum width=38mm] (close) at (0,-22mm)
    {时刻 B\\\code{close(3)}};
  \node[os memory, minimum width=43mm] (file1) at (54mm,-22mm)
    {fd 3 复用\\指向对象 G};
  \node[os compute, minimum width=48mm] (req) at (121mm,0)
    {在途请求 R\\保存对象 F 引用\\不再保存 fd 3 语义};

  \draw[os strong bus] (fd0.east) -- (file0.west);
  \draw[os strong bus] (file0.east) -- node[os label, above]{取得引用} (req.west |- file0.east);
  \draw[os feedback] (close.east) -- (file1.west);
  \draw[os guide, dashed] (file1.east) -- node[os label, below]{不能替换 R 的目标} (req.west |- file1.east);
\end{pdfdiagram}
\diagramnote{fd 是可复用槽位，不是永久身份。请求 R 在时刻 A 取得 F 的引用后，关闭与槽位复用只改变 fd table；F 要等 R 也释放引用后才能销毁。}

\subsection{取得引用和删除槽位为什么需要同步}

\code{fget(fd)} 与 \code{close(fd)} 可以并发。若读取槽位指针和增加对象引用之间没有保护，
\code{close} 可能删除最后一个槽位引用并释放对象，读取方随后对已释放内存加引用。实现
可以用锁保护这两个动作，也可以使用 RCU、原子引用和重试协议；无论采用何种技术，都必须
形成同一语义：

\begin{enumerate}
\tightlist
\item 查询要么在槽位仍有效时取得一个真实对象引用。
\item 要么观察到槽位已经关闭并返回 \code{-EBADF}。
\item 不允许返回一个已开始销毁且无法安全使用的裸指针。
\end{enumerate}

\begin{lstlisting}
lookup_fd_with_reference(fd):
  enter fd-table read-side protection
  file = read slot fd
  if file exists and reference can be acquired:
    leave protection
    return file
  leave protection
  return EBADF
\end{lstlisting}

引用计数解决“对象何时可以释放”，不自动解决“对象内部字段怎样并发修改”。取得
\code{file} 引用后，偏移、flags、队列或文件系统状态仍按各自锁和原子规则保护。生命
周期同步与内容同步是两项不同职责。

\subsection{打开对象为什么还要与 inode 分开}

两个进程可以分别打开同一路径。它们需要共享文件身份、长度和权限元数据，却可以拥有不同
当前偏移和打开 flags。若把所有状态放在 inode 中，独立打开者会相互推进偏移；若只保存
打开对象，又无法让多个打开者共享同一文件身份。因此，fd 槽位引用一次打开实例，打开
实例再引用持久文件身份。

\begin{longtable}{|L{0.18\linewidth}|L{0.28\linewidth}|L{0.40\linewidth}|}
\hline
层次 & 本章需要的边界状态 & 为什么不能由相邻层替代 \\ \hline
fd 槽位 & 槽位引用、close-on-exec 等 fd 级标志 & 整数会复用，不适合保存共享偏移或文件身份 \\ \hline
打开对象 & 当前偏移、打开模式、引用计数、操作方法 & 每次独立打开需要自己的可变访问状态 \\ \hline
inode 等对象身份 & 所有者、mode、长度和内容索引入口 & 多次打开和多个名字可以共享同一身份 \\ \hline
\end{longtable}

这里不展开文件系统布局，只建立权限检查所需的对象边界。完整 inode 为什么存在、磁盘上
保存哪些字段，将由文件系统部分从块分配问题重新推导。

\subsection{路径字符串为什么要转换成稳定路径对象}

字符串 \code{"/tmp/report"} 也不是永久对象身份。另一个线程可以在检查和打开之间
rename 路径组件，挂载关系也可能改变。内核的路径遍历因此逐级取得 dentry、mount 和
inode 等稳定引用，并在需要的位置检查目录搜索权限、符号链接策略和最终对象权限。检查
必须绑定实际解析到的对象，而不是先检查一个字符串，稍后重新解析同一字符串使用。

\begin{longtable}{|L{0.16\linewidth}|L{0.32\linewidth}|L{0.40\linewidth}|}
\hline
不安全步骤 & 并发变化 & 结果 \\ \hline
检查路径字符串对应普通文件 & 攻击者把末级名字替换成符号链接 & 使用阶段可能到达检查时未授权的目标 \\ \hline
检查父目录后释放引用 & 父目录被移动到另一挂载关系 & 后续创建发生在不同可见树中 \\ \hline
凭名字再次查找已检查对象 & 原名字删除并复用到新 inode & 检查与使用作用于两个身份 \\ \hline
\end{longtable}

\code{openat} 允许调用者用目录 fd 作为起点，减少依赖易变的当前工作目录；
\code{openat2} 一类接口还能明确限制符号链接、挂载穿越和解析范围。这些接口不是随意增加
flags，而是在基本路径解析遇到具体竞态和隔离需求后形成的约束。

\subsection{对象发布为何要区分保留槽位与安装引用}

创建新 fd 的系统调用既要建立打开对象，也要在当前 fd table 中选择空槽。若先把半初始化
对象发布到槽位，另一个线程可能立即使用未完成对象；若所有工作完成后才查找槽位，又可能
因 fd 上限失败并浪费已经取得的外部资源。常见协议把过程拆成：

\begin{enumerate}
\tightlist
\item 预留一个尚不可使用的 fd 编号。
\item 构造并验证打开对象，取得所需引用。
\item 只有对象完全可用时，原子地把对象安装到预留槽位。
\item 任一步失败都释放预留槽和临时对象。
\end{enumerate}

\begin{pdfdiagram}[scale=0.84]
  \node[os state, minimum width=34mm] (free) at (0,0)
    {空槽\\未分配};
  \node[os control, minimum width=38mm] (reserved) at (45mm,0)
    {槽位已预留\\其他查找不可用};
  \node[os memory, minimum width=38mm] (built) at (94mm,0)
    {对象已构造\\引用和权限完成};
  \node[os compute, minimum width=35mm] (visible) at (143mm,0)
    {安装完成\\fd 对用户可见};

  \draw[os strong bus] (free.east) -- node[os label, above]{reserve} (reserved.west);
  \draw[os strong bus] (reserved.east) -- node[os label, above]{build} (built.west);
  \draw[os strong bus] (built.east) -- node[os label, above]{publish} (visible.west);
  \draw[os feedback] (built.south) -- ++(0,-14mm) -|
    node[os label, below, pos=0.55]{失败：释放对象与预留槽} (free.south);
\end{pdfdiagram}
\diagramnote{“找到空编号”和“让编号引用有效对象”是两个边界。预留防止竞争分配，发布点保证其他线程只看到完整对象。}

\subsection{引用取得之后权限仍可能随状态变化}

稳定引用保证对象身份不被替换，却不保证操作永远允许。文件系统可能被重新挂成只读，对象
可能进入关闭或撤销状态，凭据可能在调用前已经改变，强制访问控制策略也可能拒绝操作。
处理函数需要在与当前操作相符的位置检查动态状态，并用对象锁、序列号或重试规则保证检查
和提交作用于同一状态版本。

这项原则可写成一个三段账本：

\begin{longtable}{|L{0.18\linewidth}|L{0.32\linewidth}|L{0.38\linewidth}|}
\hline
边界 & 必须稳定的事实 & 典型保护 \\ \hline
身份边界 & 使用期间仍是同一个打开对象、inode 或 socket & 引用计数、RCU、对象生命周期锁 \\ \hline
授权边界 & 权限判断读取一致的主体和对象安全状态 & 不可变凭据快照、对象锁、LSM hook 位置 \\ \hline
提交边界 & 状态更新与返回数量一致 & 操作锁、事务、短完成计数、错误回滚 \\ \hline
\end{longtable}

只有三项都成立，系统调用才不会出现“检查了 A、使用了 B”或“返回失败但已不可见地修改”
的边界错误。下一节继续推导授权边界中的主体状态。

\section{权限判断为何需要独立的凭据快照}

\subsection{只在 task 中保存一个 uid 为什么不够}

先设想最简单的文件权限：每个 task 只有一个用户编号，每个对象保存所有者编号和
“所有者可读”位。读取时比较两个编号即可。这个方案遇到第一个实际需求就不够了：一个
登录用户执行某个受控管理程序时，程序可能需要暂时使用文件所有者权限，完成后再恢复原
身份。系统既要记住“是谁发起”，又要记住“本次访问按谁的权限判断”。

因此，Unix 风格凭据至少区分真实身份与有效身份。真实用户 ID（real UID）描述进程来源；
有效用户 ID（effective UID）通常参与访问检查；保存的 set-user-ID 让受控程序能够在
规则允许时切换有效身份并恢复。文件访问还有 filesystem UID 一类实现字段，用于把文件
系统检查与某些临时身份变化分开。

\begin{longtable}{|L{0.20\linewidth}|L{0.28\linewidth}|L{0.40\linewidth}|}
\hline
凭据字段 & 回答的问题 & 为什么不能只保留一个 uid \\ \hline
\code{uid} & 谁启动或拥有这个进程来源 & 审计来源不应随临时授权一起消失 \\ \hline
\code{euid} & 当前访问通常按谁判断 & 受控程序需要暂时获得或放弃权限 \\ \hline
\code{suid} & 允许恢复到哪个受控身份 & 若只覆盖原值，就无法安全恢复 \\ \hline
\code{fsuid} & 文件系统检查采用哪个身份 & 让特定服务缩小文件访问身份而不改变其他语义 \\ \hline
\code{gid/egid/groups} & 当前属于哪些组 & 一个用户可能通过多个组获得对象权限 \\ \hline
\end{longtable}

字段增多不是历史装饰，而是不同语义不能安全复用一个变量。若临时提高 \code{euid} 时也
覆盖真实来源，审计记录会把请求错误归因给特权身份；若没有可恢复字段，程序只能永久放弃
或永久保留权限。

\subsection{为什么凭据要成为 task 引用的独立对象}

同一进程中的多个线程通常共享身份语义。若每个 task 都逐字段保存一份凭据，改变身份时
必须原子更新多个线程；读取方可能看见一半旧字段、一半新字段。内核更容易维护的组织是：
task 持有指向凭据对象的引用，凭据对象发布后保持不可变；身份改变时先复制旧凭据、修改
新副本、完成全部检查，再一次性替换 task 或线程组的凭据指针。

\begin{lstlisting}
credential update:
  old = current.cred
  new = copy(old)
  apply requested uid/gid/capability changes to new
  validate transition using old authority
  publish current.cred = new
  release old reference after readers leave
\end{lstlisting}

这是一种写时复制和发布协议。读取权限的热路径只需取得一个稳定凭据指针，便能在整个检查
期间看到自洽字段集合；写路径承担复制和发布成本。RCU 或等价读侧保护可延迟旧对象回收，
避免并发权限检查使用已释放凭据。

\begin{pdfdiagram}
  \node[os state, minimum width=42mm] (t1) at (0,16mm)
    {线程 T1\\\code{cred -> C0}};
  \node[os state, minimum width=42mm] (t2) at (0,-18mm)
    {线程 T2\\\code{cred -> C0}};
  \node[os memory, minimum width=48mm, minimum height=45mm] (c0) at (62mm,0)
    {不可变凭据 C0\\uid/euid/groups\\capability sets\\user namespace};
  \node[os control, minimum width=43mm] (copy) at (124mm,16mm)
    {复制并验证\\形成 C1};
  \node[os memory, minimum width=43mm] (c1) at (124mm,-18mm)
    {新凭据 C1\\一次性发布};

  \draw[os strong bus] (t1.east) -- (c0.west |- t1.east);
  \draw[os strong bus] (t2.east) -- (c0.west |- t2.east);
  \draw[os strong bus] (c0.east |- copy.west) -- (copy.west);
  \draw[os strong bus] (copy.south) -- (c1.north);
\end{pdfdiagram}
\diagramnote{权限读取使用不可变快照。身份变化不是原地逐字段修改 C0，而是构造 C1 并发布新引用；旧读者离开后 C0 才能回收。}

\subsection{当前凭据与打开时凭据为何可能不同}

处理当前系统调用时，\code{current\_cred} 描述调用者此刻的主体。某些长期对象还要记住
创建或打开时的安全上下文。例如，文件被打开后通过 fd 传给另一个进程，后者使用的打开
模式来自原打开对象，但每次操作是否还需按当前凭据检查取决于接口与安全策略。异步工作
若在内核工作线程中执行，也不能误用工作线程自身的管理身份替代原请求者。

因此，要明确两种凭据角色：

\begin{itemize}
\tightlist
\item 主体凭据：当前发起检查的 task 具有谁的权限。
\item 对象或请求凭据：对象创建、文件打开或异步请求提交时捕获的安全上下文。
\end{itemize}

异步请求若需要延后执行权限相关操作，应显式持有提交者凭据引用或已经授权的对象引用。
不能在工作线程运行时重新取 \code{current}，因为那时 \code{current} 已是内核工作线程，
它的权限通常比用户请求者更高。

\subsection{mode 位如何从一个具体访问问题得到}

已有稳定主体身份和对象身份后，最小文件检查可以把访问者分为所有者、所属组和其他主体
三类。每类分别保存读、写、执行位。执行位对目录通常表示搜索或穿越能力，而不表示读取
目录项内容；因此目录的读权限与搜索权限要分开。

\begin{longtable}{|L{0.16\linewidth}|L{0.28\linewidth}|L{0.44\linewidth}|}
\hline
请求 & 所需对象状态 & 不能替代它的检查 \\ \hline
读取普通文件内容 & 匹配主体类别的读位或等价授权 & 父目录可搜索不代表文件可读 \\ \hline
写普通文件内容 & 写位、打开模式以及非只读文件系统状态 & 拥有目录写权限不代表可改已有文件内容 \\ \hline
查找目录中的名字 & 目录搜索/执行权限 & 能列出目录名不一定能穿越到目标 \\ \hline
在目录中创建名字 & 目录写与搜索权限、挂载和策略约束 & 新文件 mode 尚未存在，不能检查它来授权创建 \\ \hline
\end{longtable}

这张表说明权限必须在拥有相应对象语义的层检查。路径遍历检查每级目录，最终打开检查目标
对象，写路径还检查打开实例和文件系统当前状态。把这些条件压成一个早期 \code{if} 会
缺少尚未解析出来的对象信息。

\subsection{ACL 为什么在三类 mode 不够时出现}

若项目文件要允许用户 Alice 读、用户 Bob 写、组 review 只读，而其他人无权访问，仅靠
所有者/组/其他三类位无法表达。改变文件所属组会影响其他组员，为每种组合创建新组也会
迅速失控。访问控制列表（access control list, ACL）为指定用户和组增加条目，并用 mask
约束一组条目的最大有效权限。

ACL 没有取消 mode 位。系统需要规定所有者条目、命名用户条目、组条目、mask 与 other
条目的求值顺序，并把 mode 与 ACL 的可见表示保持一致。权限检查因此先确定主体与对象，
再按对象是否有 ACL 选择求值规则。

\subsection{为何不能让 root 永远绕过所有检查}

单一 uid 0 能简化早期管理，却把挂载、网络配置、DAC 绕过、进程控制、模块加载等无关
能力集中到一个主体。一个只需绑定低端口的服务一旦被攻破，就可能同时获得修改文件和控制
内核的能力。由此产生的改进是把全能身份拆为若干
\term{能力}{capability}。

凭据对象通常保存多组能力位，而不只是一组“当前能力”：

\begin{longtable}{|L{0.20\linewidth}|L{0.28\linewidth}|L{0.40\linewidth}|}
\hline
能力集合 & 约束的状态 & 设计作用 \\ \hline
permitted & 进程可使其生效的上界 & 阻止任意获得不在授权上界内的能力 \\ \hline
effective & 当前权限检查实际使用的能力 & 允许暂时关闭已获准能力 \\ \hline
inheritable & 可参与 exec 继承计算的能力 & 控制普通执行链上的传递 \\ \hline
bounding & 后续 exec 也不能越过的系统性上界 & 让服务永久删去不需要的能力 \\ \hline
ambient & 满足条件时跨非特权 exec 保留的能力 & 支持不依赖 setuid 的有限授权传递 \\ \hline
\end{longtable}

能力检查必须针对具体操作选择具体 bit。\code{CAP\_DAC\_OVERRIDE} 与
\code{CAP\_NET\_ADMIN} 的职责不同；使用范围过大的能力会重新形成单一 root 的风险。
\code{CAP\_SYS\_ADMIN} 覆盖面很广，因此尤其需要避免把不相关操作都归入它。

\subsection{能力为什么必须相对 user namespace 解释}

容器内 uid 0 若等同于宿主 uid 0，容器就没有形成安全边界。user namespace 保存容器内
外 uid/gid 映射，并成为能力解释的作用域。一个凭据可以在自己的 user namespace 中拥有
某项能力，却不自动在祖先 namespace 或宿主对象上拥有同等权限。

\begin{pdfdiagram}[scale=0.88]
  \node[os state, minimum width=46mm] (cred) at (0,0)
    {容器 task 凭据\\uid 0 in U1\\effective capability};
  \node[os memory, minimum width=45mm] (u1) at (61mm,15mm)
    {user namespace U1\\映射与能力作用域};
  \node[os memory, minimum width=45mm] (u0) at (61mm,-20mm)
    {宿主 user namespace U0\\U1 的祖先};
  \node[os control, minimum width=47mm] (obj) at (127mm,0)
    {目标对象\\owner user namespace\\决定在哪一层检查};

  \draw[os strong bus] (cred.east) -- (u1.west);
  \draw[os guide] (u1.south) -- node[os label, right]{父子关系} (u0.north);
  \draw[os strong bus] (u1.east) -- (obj.west |- u1.east);
  \draw[os feedback] (u0.east) -- (obj.west |- u0.east);
\end{pdfdiagram}
\diagramnote{能力不是脱离对象的全局布尔值。检查要把调用者凭据中的能力集合放到目标对象所属 user namespace 关系中解释。}

\code{ns\_capable(target\_user\_ns, cap)} 一类语义同时回答“具有什么能力”和“相对哪个
命名空间有效”。只调用全局 \code{capable} 而忽略对象归属，可能把容器内权限扩大到宿主。

\subsection{LSM 为什么在对象已稳定后再判断}

mode、ACL 和 capability 仍主要表达 Unix 自主访问控制。系统还可能要求“Web 服务只能
读取标为 public 的对象，即使 uid/mode 原本允许”，或者“某域不能向另一域建立 socket”。
Linux Security Module（LSM）在对象语义已经明确的位置调用安全钩子，让 SELinux、
AppArmor 或 BPF LSM 等策略检查主体标签、对象标签和操作。

\begin{lstlisting}
authorize_open:
  resolve path to stable object
  check mount and object state
  evaluate mode / ACL / applicable capability
  invoke security policy hook with subject and object
  only then publish usable file reference
\end{lstlisting}

钩子过早时只有路径字符串或 fd 整数，无法绑定稳定对象；钩子过晚时对象可能已经发布或
产生副作用。合适位置是“对象已解析且引用稳定，但操作尚未提交”。不同操作的提交点不同，
所以安全钩子分布在 inode、file、task、socket、BPF 等路径，而不是只有一个总开关。

\subsection{一次授权判断的主体—对象—操作账本}

面对复杂权限问题，可以先写三项输入，再追踪决策：

\begin{longtable}{|L{0.17\linewidth}|L{0.31\linewidth}|L{0.40\linewidth}|}
\hline
输入 & 需要冻结的状态 & 代表性内容 \\ \hline
主体 & 一次检查期间一致的凭据快照 & uid/gid/groups、能力集合、user namespace、安全标签 \\ \hline
对象 & 使用期间不被替换的对象引用 & inode/socket/task、所有者、mode/ACL、对象 namespace、标签 \\ \hline
操作 & 明确的请求语义与提交点 & 读、写、执行、创建、连接、发送信号、改变配置 \\ \hline
\end{longtable}

授权结果只对这组三元关系成立。“主体能读文件 A”不能推出它能写 A，也不能推出它能读
同名但已被替换的文件 B。对象状态若在提交前发生需要重新判断的变化，操作必须在锁内检查
或使用版本号重试。

\begin{keyidea}
凭据之所以成为独立对象，是为了让一次权限检查读取自洽且可共享的主体快照。uid、有效
身份、组和多组 capability 分别表达不同授权边界；user namespace 决定能力作用域；
mode、ACL 与 LSM 按已稳定对象的语义继续收缩权限。授权从来不是入口处一个孤立的
\code{if}，而是主体、对象和操作在提交前形成的关系。
\end{keyidea}

\section{阻塞以后如何保留同一项系统调用}

\subsection{没有数据时不能靠入口帧一直占用 CPU}

现在考虑 \code{read} 一个空管道。对象引用有效，用户缓冲区范围也可写，但管道缓冲区
暂时没有数据。让内核循环检查会持续占用 CPU；立即返回错误又会迫使用户程序频繁重试。
系统需要把“当前条件不满足”保存到等待关系中，让其他 task 运行；数据到达时再使读者具备
运行资格。

入口帧只能保存用户返回现场，不能回答“等待哪个条件”。管道对象因此维护等待队列；读
task 在栈上的等待记录或长期等待对象中保存队列链接、task 引用、唤醒规则和等待原因。

\begin{longtable}{|L{0.20\linewidth}|L{0.28\linewidth}|L{0.40\linewidth}|}
\hline
等待状态 & 保存在哪里 & 作用 \\ \hline
用户返回现场 & task 内核栈入口帧 & 系统调用最终完成后回到用户程序 \\ \hline
内核继续位置 & 内核调用链与调度切换帧 & task 再次运行时从条件复查处继续 \\ \hline
等待条件关系 & 管道等待队列条目 & 数据到达路径知道应唤醒哪些 task \\ \hline
任务可运行状态 & task 与 runqueue & 调度器知道它当前是否有资格竞争 CPU \\ \hline
已完成前缀 & 处理函数局部或请求对象 & 信号、错误发生时决定返回数量还是负错误 \\ \hline
\end{longtable}

\subsection{为什么入队和检查条件必须形成原子观察边界}

若 task 先检查“管道为空”，尚未加入等待队列时写者放入数据并执行唤醒，唤醒看不到任何
等待者；读者随后入队并睡眠，可能永远等不到已经发生的事件。这是丢失唤醒。

正确协议在对象锁或等价同步下安排状态：

\begin{lstlisting}
read_from_pipe:
  lock pipe
  while pipe is empty:
    add current task to pipe.read_waiters
    set task state to interruptible sleep
    unlock pipe
    schedule
    lock pipe
    remove wait entry if still linked
  copy available data
  set task state to running
  unlock pipe
\end{lstlisting}

具体内核通常把这些步骤封装在等待原语中，并仔细安排内存顺序。核心不变量是：写者要么
观察到等待者并唤醒，要么读者在决定睡眠前观察到数据已经到达。不存在“事件发生了，但两方
都没观察到”的间隙。

\begin{pdfdiagram}
  \node[os state, minimum width=46mm] (reader) at (0,20mm)
    {读 task T\\条件：管道非空};
  \node[os memory, minimum width=46mm] (waitq) at (61mm,20mm)
    {管道等待队列\\条目指向 T};
  \node[os control, minimum width=46mm] (writer) at (122mm,20mm)
    {写 task U\\加入数据并唤醒};

  \node[os state, minimum width=46mm] (sleep) at (0,-23mm)
    {T：sleeping\\不在 runqueue};
  \node[os compute, minimum width=46mm] (ready) at (61mm,-23mm)
    {条件已成立\\从 wait queue 移除};
  \node[os state, minimum width=46mm] (runq) at (122mm,-23mm)
    {T：runnable\\进入某 CPU runqueue};

  \draw[os strong bus] (reader.east) -- node[os label, above]{锁内登记} (waitq.west);
  \draw[os strong bus] (writer.west) -- node[os label, above]{锁内发布数据} (waitq.east);
  \draw[os strong bus] (reader.south) -- (sleep.north);
  \draw[os strong bus] (waitq.south) -- (ready.north);
  \draw[os strong bus] (ready.east) -- (runq.west);
\end{pdfdiagram}
\diagramnote{等待队列保存条件关系，runqueue 保存运行资格。唤醒只把 T 变为 runnable，不保证 T 立即获得 CPU；恢复后 T 仍要在锁内复查条件。}

\subsection{唤醒为什么不等于系统调用已经完成}

数据到达时，写者只能证明“值得让读者重新检查”。数据可能被另一个读者先取走，信号也
可能先到达；因此 T 获得 CPU 后必须重新检查管道状态。等待循环使用条件而不是单次
\code{if}，既处理竞争，也处理无条件或虚假唤醒。

唤醒路径不直接改写 T 的用户返回值。它改变 task 的运行资格；T 恢复内核调用链后，自己
复制数据、推进对象状态并写入入口帧的结果槽。这样，完成数量仍由持有操作上下文的执行流
决定。

\subsection{信号为什么能打断可中断等待}

用户可能要求终止或处理某个事件。若所有阻塞调用都不可中断，task 可能长期无法响应。
可中断等待把“数据到达”和“有待处理信号”都作为离开睡眠的条件。task 被信号唤醒后，
处理函数先看是否已有提交前缀：

\begin{longtable}{|L{0.20\linewidth}|L{0.28\linewidth}|L{0.40\linewidth}|}
\hline
调用进度 & 信号到达后的优先结果 & 原因 \\ \hline
尚未提交任何数据 & 返回中断语义或准备重启 & 调用者尚不需要处理成功前缀 \\ \hline
已经读取 600 字节 & 返回 600 & 完成数量比中断错误更重要，避免重复消费 \\ \hline
操作不可安全重启 & 返回 \code{-EINTR} 等接口结果 & 从头执行可能重复产生副作用 \\ \hline
操作可重启且策略允许 & 内部标记重启，信号处理后恢复请求 & 向用户隐藏无副作用的中断窗口 \\ \hline
\end{longtable}

这里再次出现“结果优先级”。同一次调用内部可能同时发生信号和部分完成，ABI 只能返回一个
主结果。接口必须选择能让调用者正确继续的那个结果，并把信号保留到返回工作路径处理。

\subsection{系统调用重启为什么不只是把 RIP 后退}

无参数、无副作用的简单等待或许能重新执行入口指令。带超时、临时内核状态或已规范化参数
的调用还要保存剩余时间和重启函数。restart block 一类结构可以记录：

\begin{lstlisting}
restart state:
  restart_function
  normalized arguments
  remaining timeout or absolute deadline
  object identity or restart cookie
\end{lstlisting}

采用绝对 deadline 比反复使用原相对超时更可靠。若每次信号后都重新等待完整 10 秒，多次
信号会无限延长总等待；保存绝对截止时间后，重启只计算剩余预算。

\begin{lstlisting}
restart_timed_wait:
  remaining = deadline - current_time
  if remaining <= 0:
    return timeout
  resume wait with remaining budget
\end{lstlisting}

并非所有调用都允许自动重启。创建对象、发送消息或设备控制若已产生不可辨认的副作用，
重新执行可能重复操作。重启规则属于每项系统调用的语义，不由入口层统一猜测。

\subsection{调度切换后按什么顺序继续}

读 task T 睡眠后，调度器保存 T 的内核切换现场并运行 U。数据到达使 T runnable；当
调度器后来选择 T 时，恢复顺序是：

\begin{enumerate}
\tightlist
\item 恢复 T 的地址空间与体系结构线程状态（若此前运行的是另一进程）。
\item 恢复 T 的内核栈指针和被调用者保存寄存器。
\item 从 \code{schedule} 之后的内核位置继续，而不是直接回用户态。
\item 重新取得对象锁并复查等待条件。
\item 形成数据或错误结果，逐层返回到系统调用出口。
\item 出口处理信号、调度请求和用户返回状态，最后恢复入口帧。
\end{enumerate}

入口帧一直位于 T 的内核栈中，但调度器首次恢复的不是该帧，而是更靠近当前栈顶的切换
现场。只有系统调用 C 调用链完成后，控制流才逐步退回入口帧。

\begin{keyidea}
阻塞不会把系统调用变成另一项请求。入口帧保存用户继续状态，内核调用链保存处理函数继续
位置，等待队列保存条件关系，runqueue 保存运行资格，局部进度保存已提交前缀。唤醒只改变
运行资格；task 恢复后必须重新检查条件并亲自形成返回结果。
\end{keyidea}

\section{进入内核以后为何仍要区分执行上下文}

\subsection{系统调用执行期间设备事件不会停止}

task 进入内核后，网卡仍可能收到数据，定时器仍会到期，其他 CPU 也仍在修改共享对象。
若整个系统调用期间都禁止中断，一次缓慢的路径解析或页面回收就会让设备完成长时间得不到
处理，调度时钟也无法推进。若始终允许中断，又会出现新的问题：同一 CPU 正在使用 task 的
内核栈和每 CPU 临时状态时，中断入口会再次保存一份现场。

于是不能只说“当前在内核态”。至少要知道控制流由哪一种事件进入、能否阻塞、能否被抢占，
以及它正在使用哪一份栈。系统调用上下文与中断上下文都运行内核代码，但它们拥有的继续
条件不同：

\begin{longtable}{|L{0.20\linewidth}|L{0.28\linewidth}|L{0.40\linewidth}|}
\hline
执行来源 & 当前控制流依附的状态 & 能否主动等待 \\ \hline
系统调用 & 当前 task、其内核栈和入口帧 & 在不持有禁止睡眠资源时可以 \\ \hline
用户异常 & 当前 task、异常帧和故障地址 & 可修复故障路径可能等待页面或 I/O \\ \hline
设备中断 & CPU 上被打断的现场与中断入口帧 & 不能把当前 task 当成中断工作的请求者去睡眠 \\ \hline
不可屏蔽事件 & 最小独立入口状态与严格受限的代码 & 不能依赖普通锁或普通调度 \\ \hline
\end{longtable}

设备中断恰好发生在 task T 运行时，不表示中断工作属于 T。T 只是被借用了 CPU；中断
完成后通常回到被打断位置。若中断处理因等待而调用调度器，就没有一项独立 task 可保存
“中断处理到哪里”，也会把 T 的内核调用链夹在一个无法解释的等待关系中。

\subsection{嵌套入口为什么需要另一份入口帧}

设 T 正在执行系统调用处理函数，处理器在位置 \(K_1\) 收到网卡中断。原系统调用入口帧
保存的是“将来怎样回到用户位置 \(U_0\)”；中断入口新保存的则是“怎样回到内核位置
\(K_1\)”。两份帧不能合并：

\begin{center}
\begin{tikzpicture}[
  >=Latex,
  box/.style={draw=BookRule, rounded corners=2pt, align=left,
              minimum width=9.5cm, minimum height=0.74cm, fill=BookBlueTint},
  arr/.style={->, thick, draw=BookAccent}
]
\node[box] (u) at (0,3.8) {较低地址：系统调用入口帧——用户 \code{RIP=U0}、用户 \code{RSP}、原参数};
\node[box] (call) at (0,2.55) {系统调用内核调用链——对象引用、局部进度、返回地址};
\node[box] (irq) at (0,1.30) {中断入口帧——内核 \code{RIP=K1}、当时标志、被打断寄存器};
\node[box] (handler) at (0,0.05) {较高地址：中断处理调用链——设备状态与完成记录};
\draw[arr] (u) -- (call);
\draw[arr] (call) -- (irq);
\draw[arr] (irq) -- (handler);
\end{tikzpicture}
\end{center}
\diagramnote{中断返回只弹出中断调用链和中断帧，恢复到 \(K_1\)。系统调用稍后完成时才使用
更早的系统调用入口帧回到 \(U_0\)。}

若中断发生在用户态，中断帧保存用户继续状态；若发生在内核态，它保存内核继续状态。入口
代码必须依据被打断特权级和入口类型决定需要切换哪一份栈、需要补齐哪些字段，而不能假设
所有事件都来自用户态。

\subsection{为什么中断入口常使用每 CPU 或专用栈状态}

中断在任意 task 上都可能发生。若它先查找一个需要取得普通锁的全局结构来决定栈位置，
而中断恰好打断了持有该锁的代码，入口尚未建立安全现场就会自锁。最早阶段更适合使用
当前 CPU 可直接定位的状态，包括入口栈顶、嵌套深度、当前 task 指针和紧急事件专用栈。

专用栈还处理两类风险。第一，当前 task 的内核栈可能接近耗尽，严重故障若继续压栈会覆盖
邻接内存；第二，某些异常可能在普通入口本身故障时到达，需要与已经损坏的栈状态隔离。
专用栈不是为了让所有中断拥有独立线程，而是确保最早保存与诊断不依赖可能失效的普通栈。

\begin{longtable}{|L{0.22\linewidth}|L{0.29\linewidth}|L{0.37\linewidth}|}
\hline
状态位置 & 适合保存的内容 & 不应混入的内容 \\ \hline
task 内核栈 & 本次系统调用调用链、局部进度、入口帧 & 其他 task 的入口现场 \\ \hline
每 CPU 入口区 & 最早栈指针、当前 task、入口嵌套标志 & 会跨 CPU 迁移的长期请求状态 \\ \hline
普通中断栈 & 可屏蔽中断的短调用链与临时设备状态 & 需要睡眠等待的操作进度 \\ \hline
紧急专用栈 & 严重异常的最小保存与诊断状态 & 常规业务缓冲区和大对象 \\ \hline
\end{longtable}

CPU 迁移进一步限制了每 CPU 状态的使用。T 读取“当前 CPU 的临时槽”后若允许抢占并迁移，
后续写回可能落到另一 CPU。访问方要么在短临界区禁止迁移，要么保存并使用显式 CPU 身份，
不能把一次读取的 CPU 编号当作永久 task 属性。

\subsection{为何需要记录当前上下文允许哪些动作}

一段公共辅助函数可能从系统调用路径调用，也可能从设备中断调用。仅凭代码地址无法判断
它能否睡眠。内核需要维护足够的上下文状态，例如中断嵌套、抢占禁止深度、持有的自旋类
资源以及中断屏蔽状态，使调度与诊断路径能够拒绝不合法的等待。

\begin{lstlisting}
may_block_here:
  require running on behalf of a schedulable task
  require not inside hard interrupt or emergency entry
  require preemption/atomic nesting permits scheduling
  require no held resource whose protocol forbids sleep
\end{lstlisting}

“原子上下文”在这里不是说整段代码由一条硬件原子指令完成，而是说当前执行不能把 CPU
让给另一个 task 后再继续。原因可能是处于中断调用链，也可能是持有自旋锁或显式禁止抢占。
把这两个含义混为“执行很快”会造成错误：一段代码即使只有几十条指令，只要调用了可能缺页
或等待内存的函数，就可能睡眠。

\subsection{用户复制为何会影响锁的选择}

\code{copy\_from\_user} 看似只是复制字节，但目标用户页可能尚未建立 PTE，需要进入缺页
路径；缺页处理又可能分配页面或等待存储 I/O。因此用户复制在一般情形下是潜在睡眠点。

若对象实现持有禁止睡眠的自旋锁再复制用户内存，可能出现：

\begin{enumerate}
\tightlist
\item 复制访问用户页并产生 page fault；
\item fault 路径需要等待页面；
\item 调度器发现当前处于禁止调度的上下文；
\item 操作无法安全继续，也不能保留自旋锁让其他 CPU 无限等待。
\end{enumerate}

正确安排要由对象协议决定。常见方向包括先把小输入复制到内核快照，再取得对象自旋锁；
或者在可睡眠互斥边界内复制；或者为大 I/O 预先固定并验证页面，然后在不可睡眠阶段只访问
已保证存在的映射。三者的代价和生命期不同，但都必须明确“缺页可能发生在哪个边界”。

\begin{longtable}{|L{0.21\linewidth}|L{0.28\linewidth}|L{0.39\linewidth}|}
\hline
安排方式 & 获得的保证 & 仍需处理的问题 \\ \hline
先复制后加锁 & 持锁阶段不再访问原用户输入 & 加锁前对象状态可能变化，锁内必须重验 \\ \hline
可睡眠锁内复制 & 对象条件在复制期间保持稳定 & 长缺页会延长锁占用，影响其他 task \\ \hline
预先固定用户页 & 后续阶段可引用稳定页面集合 & 固定成本、解除固定和写回责任 \\ \hline
分段保留并复制 & 限制单次锁占用与故障范围 & 必须精确推进已提交前缀 \\ \hline
\end{longtable}

这也解释了为何“检查参数—加锁—复制—提交”不是所有调用都能照搬的固定顺序。顺序取决于
每一步可能睡眠、依赖哪些对象状态，以及失败时能否撤销。

\subsection{中断工作为什么要拆成短入口和可调度后续}

设备完成可能需要解析大量描述符、分配缓冲区、唤醒 task 或继续网络协议处理。全部放在
硬中断上下文会延长同 CPU 上其他中断的响应时间，也不能使用可睡眠操作。于是先在严格
入口中完成不可延迟的最小工作：确认设备事件、取得完成描述符的所有权、屏蔽或应答必要
状态，并登记后续工作。

后续可在受限制的延后上下文或内核线程中执行。若工作最终可能等待，就需要一项可调度的
task 作为载体。拆分后的所有权边界必须明确：入口把哪些描述符移交给后续队列，设备何时
可以复用槽位，后续失败由谁释放缓冲区。

\begin{center}
\begin{tikzpicture}[
  >=Latex,
  node distance=9mm and 12mm,
  box/.style={draw=BookRule, rounded corners=2pt, align=center,
              minimum width=3.25cm, minimum height=1.0cm, fill=BookBlueTint},
  arr/.style={->, thick, draw=BookAccent}
]
\node[box] (dev) {设备置完成状态\\触发中断};
\node[box, right=of dev] (hard) {短入口\\确认并取得描述符};
\node[box, right=of hard] (defer) {后续队列\\拥有完成状态};
\node[box, below=13mm of defer] (task) {可调度执行\\协议处理或唤醒};
\node[box, left=of task] (resume) {原系统调用 task\\复查等待条件};
\draw[arr] (dev) -- (hard);
\draw[arr] (hard) -- (defer);
\draw[arr] (defer) -- (task);
\draw[arr] (task) -- (resume);
\end{tikzpicture}
\end{center}
\diagramnote{设备中断不直接替等待中的 task 完成系统调用。它取得硬件完成状态并安排后续；
后续更新对象条件、唤醒 task，task 恢复后再形成自己的系统调用结果。}

\subsection{内核抢占为何不同于中断嵌套}

中断是外部或处理器事件插入当前控制流；内核抢占则是调度器在安全点暂停当前 task，选择
另一个 task。两者都会让当前代码暂时不执行，但保存对象与恢复位置不同。

\begin{longtable}{|L{0.20\linewidth}|L{0.30\linewidth}|L{0.38\linewidth}|}
\hline
暂停原因 & 主要保存位置 & 恢复关系 \\ \hline
设备中断插入 & 当前栈上的中断帧或专用中断栈 & 中断返回到被打断指令 \\ \hline
内核抢占 T & T 的内核栈与调度切换帧 & 调度器将来选择 T 后继续 \\ \hline
系统调用主动等待 & T 的调用链、等待队列与切换帧 & 条件唤醒后由调度器恢复 \\ \hline
用户态信号交付 & 用户信号帧与改写后的入口帧 & 先运行处理函数，信号返回后恢复原现场 \\ \hline
\end{longtable}

代码在修改每 CPU 数据或持有某些低层锁时，可以短暂禁止内核抢占，保证 task 不会在临界
区被迁移或切换。禁止抢占不一定屏蔽设备中断；如果中断处理也访问同一状态，还需要相应的
中断同步协议。反过来，屏蔽本地中断也不自动构成跨 CPU 互斥，其他 CPU 仍可并发访问。

\subsection{怎样核对一段入口代码的上下文闭合}

对入口到处理函数之间的每段代码，可按以下状态逐项核对，而不是只看函数是否“很短”：

\begin{enumerate}
\item 当前使用 task 栈、中断栈还是紧急栈？最大嵌套深度怎样受限？
\item 当前代码代表哪个 task，还是不代表任何可调度 task？
\item 本地中断是否允许，内核抢占是否允许，task 是否可能迁移 CPU？
\item 持有哪些锁或引用；其中哪些要求释放前不能 fault、不能调度？
\item 访问用户地址是否可能缺页；内存分配是否可能回收或等待？
\item 事件再次嵌套时，会覆盖每 CPU 临时槽还是获得独立槽位？
\item 离开时哪一份帧被弹出，返回内核位置还是用户位置？
\end{enumerate}

这些问题将“上下文错误”还原为具体状态冲突。例如，“在中断中睡眠”实质是没有一项可独立
调度的 task 来拥有等待中的调用链；“持自旋锁访问用户页”实质是故障路径可能要求调度，而
当前锁协议禁止让出 CPU；“每 CPU 指针失效”实质是读取与使用之间发生了迁移。

\begin{keyidea}
内核态不是单一执行环境。系统调用、用户异常、设备中断、紧急事件和调度恢复拥有不同的栈、
继续位置与阻塞能力。每段代码只有在说明了当前现场由谁拥有、允许哪些嵌套以及能否调度后，
才能判断锁、用户复制和内存分配是否安全。
\end{keyidea}

\section{边界观察为何也必须服从状态形成顺序}

\subsection{只记录入口参数为何无法说明实际操作}

诊断程序看到 \code{openat(dirfd, user\_path, flags)} 的入口时，只能得到整数、标志位和一个
用户地址。路径内容尚未形成内核快照，\code{dirfd} 尚未转换成稳定目录对象，符号链接和
挂载关系也尚未解析。若日志此时写下“打开了某文件”，它把请求意图误写成已经发生的事实。

反过来，只在出口记录返回值也不够。一个 \code{-EACCES} 不能说明拒绝的是哪个对象、使用
哪份凭据、在哪一层策略中发生。可核对的边界记录应随着信息稳定逐步补全：

\begin{longtable}{|L{0.19\linewidth}|L{0.28\linewidth}|L{0.41\linewidth}|}
\hline
观察时刻 & 此时稳定的信息 & 不能提前声称的事实 \\ \hline
入口建立后 & ABI、调用编号、原始标量参数 & 用户指针内容真实且不变 \\ \hline
参数快照后 & 内核拥有的字符串或结构字段 & 路径已对应最终对象 \\ \hline
对象解析后 & 稳定对象身份、挂载与命名空间关系 & 权限已经允许、操作已经提交 \\ \hline
授权完成后 & 主体、对象、操作与策略判断 & 后续对象状态仍允许提交 \\ \hline
提交边界后 & 实际改变、完成前缀与错误 & 数据已经持久化或被外部系统接收 \\ \hline
出口工作后 & 最终返回值、信号或重启决定 & 用户程序已经处理该结果 \\ \hline
\end{longtable}

同一条审计事件可以在入口创建上下文，随后逐步填充字段，最后在出口发布。若调用阻塞，
上下文必须跟随 task 或独立请求状态保存，不能留在某个 CPU 的临时槽中；task 恢复时可能
已经迁移到另一 CPU。

\subsection{原始参数与规范化参数为什么都可能需要保留}

原始参数回答“调用者提出了什么”，规范化参数回答“内核按哪种含义执行”。两者可能不同。
例如旧 ABI 的窄时间字段会转换为统一宽类型；路径中的相对起点会结合目录对象；标志组合会
被检查并分解成内部操作选项。

安全日志不能随意重新读取用户指针来补字段，因为用户已经可以改写原内存。若确实要记录
字符串，应使用本次操作已经取得的内核快照，并设置长度上限与转义规则。否则日志本身会
产生第二次 TOCTOU：系统实际检查的是字符串 A，日志稍后重新读取成字符串 B。

\begin{longtable}{|L{0.22\linewidth}|L{0.28\linewidth}|L{0.38\linewidth}|}
\hline
记录形式 & 适合回答的问题 & 必须注明的限制 \\ \hline
原始寄存器值 & 调用者如何编码请求 & 指针只是地址，尚无对象语义 \\ \hline
内核参数快照 & 实现实际解析了哪些字段 & 可能已经过兼容转换和规范化 \\ \hline
稳定对象标识 & 操作最终指向哪个内核对象 & 路径名称后来可以改变 \\ \hline
提交计数与终态 & 实际产生多少效果 & 完成不一定代表持久化 \\ \hline
\end{longtable}

\subsection{跟踪器能够改写入口帧时谁拥有最终请求}

某些跟踪机制允许监管者在入口暂停 task，查看甚至修改系统调用号、参数寄存器或返回值。
暂停发生时，被跟踪 task 的入口帧已经存在，但对象操作尚未开始。监管者完成修改后，内核
必须重新读取入口帧并按新值执行分发与验证，不能继续使用暂停前缓存的编号。

这形成明确的修改边界：

\begin{enumerate}
\tightlist
\item task 建立入口帧，产生原始请求状态；
\item 跟踪入口事件把 task 停在对象副作用之前；
\item 监管者通过受控接口修改允许字段；
\item task 恢复，边界层把修改后的帧视为最终原始请求；
\item 参数快照、对象解析和授权从该请求重新开始。
\end{enumerate}

监管者不能因为有能力观察入口，就直接获得任意内核指针或绕过用户地址检查。它改的是被
跟踪 task 的用户可见请求状态；修改后的指针仍然属于该 task 的地址空间，修改后的操作仍
要经过对象权限与安全策略。

出口跟踪同样位于对象结果已经形成之后、处理器返回用户态之前。监管者可以观察规定的结果
槽；若接口允许改写返回值，审计记录应能区分“对象实现结果”和“最终暴露结果”。改变返回
数字并不会自动撤销已经提交的对象副作用。

\subsection{入口策略拒绝为什么要发生在对象副作用之前}

系统调用集合过滤器适合根据编号、体系结构和原始标量参数快速收缩请求面。若它允许通知
一个用户态监管者再决定，原 task 可能等待较长时间。等待期间用户内存、fd 槽和进程关系
都可能改变，所以监管者不能仅凭入口时看到的 fd 数字替内核完成最终对象授权。

安全分工是：入口策略决定“这类请求是否可以继续进入内核实现”；对象实现随后取得稳定引用
并判断“当前主体能否对当前对象执行该操作”。若监管协议要代替某项操作，它必须定义参数
快照、返回值注入、fd 等对象如何安全移交，以及原 task 退出时怎样取消通知。否则一次
入口批准会被错误地当作对未来任意 fd 重用的批准。

\subsection{阻塞调用的时长为何要拆成几个区间}

从入口到出口的总时长包含不同原因。只报告一个总数，无法区分边界开销、锁竞争、条件等待
和设备服务。可以使用已经得到的状态边界划分时间：

\begin{longtable}{|L{0.23\linewidth}|L{0.29\linewidth}|L{0.36\linewidth}|}
\hline
时间区间 & 起止边界 & 主要解释对象 \\ \hline
入口固定时间 & 入口指令到对象实现开始 & 现场保存、过滤、分发与参数转换 \\ \hline
对象活动时间 & 获得对象到首次等待或完成 & 锁、查找、复制与提交算法 \\ \hline
非运行等待时间 & 离开 runqueue 到重新获得 CPU & 条件等待与调度延迟 \\ \hline
唤醒后排队时间 & 变为 runnable 到实际恢复 & 运行队列竞争与优先级 \\ \hline
出口固定时间 & 结果确定到用户态恢复 & 信号、跟踪、调度检查与返回验证 \\ \hline
\end{longtable}

时间源本身也要注明。墙上时间可能被管理员调整，单调时钟适合计算持续时间，CPU 运行时间
只累计 task 真正执行的区间。把不同来源的时间戳直接相减，可能得到负值或把睡眠误算为
处理器工作。

\subsection{观察代码为何不能破坏原路径}

跟踪点位于锁内、入口最早阶段或中断上下文时，记录动作受同样的上下文限制。它不能无界
分配内存、等待用户态消费者或复制任意长字符串。常见安排是写入预分配的每 CPU 环形缓冲，
记录固定大小的标量和稳定标识，随后由可调度读取者解释。

即使这样仍要定义缓冲区满时的策略：覆盖旧记录、丢弃新记录或增加丢失计数。不能让关键
中断路径为保住日志而无限等待。日志记录还要避免泄露未初始化填充字节、内核地址和不属于
观察者权限范围的数据。

观察设施关闭时，热路径应尽量只付出一次可预测的关闭检查；启用后增加的成本要能被测量。
这属于正确路径建立后的优化。不能为减少跟踪开销而让入口记录与出口记录失去配对，也不能
因审计分配失败就悄悄执行一项本应被强制审计策略拒绝的敏感操作。

\begin{keyidea}
观察不是站在系统调用之外读取一条现成事实。请求、参数快照、对象身份、授权结果、提交
前缀和最终用户结果在不同时刻形成。日志、跟踪和性能计时只有绑定这些边界，才能说明记录
的是意图、稳定输入、实际对象还是已经提交的效果。
\end{keyidea}

\section{处理函数结束以后为何还不能立即返回}

\subsection{一个返回值不足以描述返回时刻的全部状态}

设管道读取已经得到 600 字节，处理函数把 600 写入入口帧的结果槽。若内核此时只恢复
\code{RIP}、\code{RSP} 和通用寄存器，表面上已经完成了调用。然而在 task 阻塞期间，
另外几类状态可能同时发生变化：

\begin{itemize}
\tightlist
\item 时间片已经用完，task 应先让出处理器；
\item 一个信号已经挂起，需要在用户态继续执行前交付；
\item 跟踪器要求观察或改写本次调用的退出状态；
\item 审计设施需要把最终结果与入口记录配对；
\item 调用被内部标记为可重启，必须决定返回错误还是安排重试；
\item 入口帧中的用户返回地址或标志不再满足快速返回指令的约束。
\end{itemize}

这些状态不属于管道读取算法，却决定 task 下一步实际执行什么。若每项系统调用都自行处理，
不同调用会产生不同的信号、调度和跟踪语义。于是需要一段共享的返回工作路径：对象实现只
提交操作结果；返回路径在结果已经确定后，统一结清 task 级的未决工作。

\begin{longtable}{|L{0.20\linewidth}|L{0.27\linewidth}|L{0.41\linewidth}|}
\hline
未决状态 & 返回前要做的动作 & 不能交给对象实现的原因 \\ \hline
需要重新调度 & 调用调度器，之后重新检查标志 & 时间片属于 task 与 CPU，不属于文件或管道 \\ \hline
挂起信号 & 选择信号、建立用户信号帧 & 信号会改写用户继续位置，与具体对象无关 \\ \hline
调用重启 & 改写结果与继续状态，保留重启参数 & 是否重启取决于信号处置和已提交前缀 \\ \hline
跟踪或审计 & 生成退出事件，允许受控观察 & 入口与退出必须共享统一边界 \\ \hline
快速返回不安全 & 改走能完整恢复状态的慢路径 & 返回指令受体系结构条件限制 \\ \hline
\end{longtable}

返回工作不是一次性检查。处理信号时可能唤醒调试器，恢复后时间片又可能耗尽；调度回来后
还可能收到新信号。合理结构是“观察工作标志—完成一种工作—重新观察”，直到没有必须在
用户态之前完成的事项。

\begin{lstlisting}
leave_system_call:
  repeat:
    flags = task.return_work
    if flags says reschedule:
      schedule()
      continue
    if flags says signal or restart:
      prepare_user_delivery()
      continue
    if flags says trace or audit:
      finish_boundary_observation()
      continue
  until no return work remains
  validate_user_return_state()
  restore_and_return()
\end{lstlisting}

这里的循环不是无条件忙等。每个分支都消费一个状态或发生一次有意义的切换；重新读取标志
是为了覆盖并发到达的新事件。

\subsection{为什么调度检查位于用户返回之前}

系统调用处理函数结束时，task 正在内核态运行。若时间片已经耗尽但仍直接返回用户态，
task 会继续运行，直到下一次中断或入口才可能让出 CPU。负载很高时，这会破坏调度延迟的
上界。

返回边界已经具备三个条件：当前 task 有完整可恢复现场；内核栈和地址空间仍处于可控状态；
处理函数不再持有要求立即释放的局部锁。因此它是安全调度点。调度器保存的仍是内核切换
现场。将来恢复该 task 时，它从返回工作循环继续，重新读取标志，而不是绕过未处理的信号
直接进入用户态。

还要区分“需要调度”和“已经选择了下一个 task”。前者只是 task 或 CPU 上的请求标志；
只有调用调度器并取得运行队列锁后，才会做出选择。若此时没有更合适的可运行任务，当前
task 仍可能继续执行。

\subsection{信号交付为何要在用户栈上建立一份新现场}

信号处理函数是用户代码。内核不能在自己的栈上调用它，也不能丢失原来的用户执行状态。
要让处理函数结束后回到原位置，必须在用户地址空间中建立一份可由信号返回协议解释的
现场。它通常包含：

\begin{longtable}{|L{0.22\linewidth}|L{0.27\linewidth}|L{0.39\linewidth}|}
\hline
信号帧内容 & 来源 & 恢复时的用途 \\ \hline
被中断的通用寄存器 & 系统调用入口帧或用户异常帧 & 让原用户控制流继续 \\ \hline
原用户 \code{RIP/RSP/flags} & 已验证的入口状态 & 恢复指令位置、栈和允许的标志 \\ \hline
信号编号与扩展信息 & 内核挂起信号对象 & 作为处理函数参数 \\ \hline
信号屏蔽集合 & task 的信号状态 & 处理函数结束后恢复屏蔽关系 \\ \hline
浮点与向量状态 & task 的体系结构扩展状态 & 避免处理函数破坏原计算 \\ \hline
用户返回跳板地址 & 受信任的用户映射或 ABI 约定 & 发起受控的信号返回请求 \\ \hline
\end{longtable}

这份帧写在用户栈或预先登记的备用信号栈上，所以仍要走用户内存验证和受控复制。普通用户
栈可能已经接近边界，备用栈可能未映射，处理器扩展状态也可能使帧比预期更大。构造失败
不能退回到一个半写的帧上继续执行；内核需要把失败转换为规定的致命信号或终止结果。

\begin{center}
\begin{tikzpicture}[
  >=Latex,
  box/.style={draw=BookRule, rounded corners=2pt, align=left,
              minimum width=9.2cm, minimum height=0.72cm, fill=BookBlueTint},
  arr/.style={->, thick, draw=BookAccent}
]
\node[box] (old) at (0,3.3) {原用户栈：调用者局部变量与返回地址};
\node[box] (frame) at (0,2.0) {新信号帧：原寄存器、屏蔽集合、扩展状态、返回跳板};
\node[box] (handler) at (0,0.7) {返回后的用户现场：\code{RIP=handler}，\code{RSP=signal frame}};
\draw[arr] (old) -- node[right, font=\small]{向下预留并验证空间} (frame);
\draw[arr] (frame) -- node[right, font=\small]{改写入口帧中的用户继续状态} (handler);
\end{tikzpicture}
\end{center}
\diagramnote{内核并不是在用户栈上保存自己的调用链。信号帧只保存用户可恢复现场；系统调用
内核栈仍按正常出口退栈。}

信号处理函数最终执行的“返回”也不能只是普通 \code{ret}。它需要请求内核读取信号帧，
恢复屏蔽集合和扩展状态，并逐字段验证用户可控制的数据。信号帧位于用户内存，用户可能在
处理期间修改它；恢复路径必须把它当作不可信输入，不能因为它最初由内核写入就跳过检查。

\subsection{重启决定为什么必须晚于信号选择}

内核在等待中断时可能暂存“这项调用可以重启”。这还不是最终决定。若用户安装了自定义
处理函数，内核需要先构造信号帧；若信号被忽略，可能不必改变用户控制流；若调用已经提交
部分结果，则应优先返回完成数量。

最终重启至少要同时满足四项条件：

\begin{enumerate}
\tightlist
\item 当前调用的语义允许从保存的状态继续或重新进入；
\item 尚无必须先向用户报告的提交前缀；
\item 所选信号的处置方式允许自动重启；
\item 重启参数、剩余期限与对象身份仍然有效。
\end{enumerate}

满足条件时，返回路径可以把入口帧改造成“信号处理结束后再次提出请求”的状态。这里要
恢复原系统调用编号，而不是把当前结果寄存器中的内部重启码当成新编号。若使用 restart
block，则信号返回后进入专门的重启入口，再由该入口使用保存的规范化参数。

\subsection{为什么返回地址和标志必须再次验证}

入口时保存的用户 \code{RIP} 与 \code{RSP} 可能被信号、跟踪器或系统调用语义受控修改。
它们也可能来自恶意用户构造的信号返回帧。内核即将把执行权交还给较低特权级，此时必须
确认：

\begin{itemize}
\tightlist
\item 指令地址和栈地址具有体系结构规定的用户地址形式；
\item 返回目标位于允许的用户地址范围，而不是内核映射；
\item 段选择子、特权级和地址宽度与当前 ABI 一致；
\item 标志寄存器中的特权控制位已清除或恢复为规定值；
\item 调试、单步、地址空间和控制流保护相关状态走了相应慢路径；
\item 不会把仅用于内核的临时寄存器内容泄露给用户态。
\end{itemize}

“入口时验证过”不能代替此处验证，因为中间路径可能合法地改变返回状态。返回验证保护的
不是系统调用参数，而是下一条用户指令的执行边界。

\subsection{快速返回为何只能是一项经过证明的优化}

x86-64 提供快速系统调用返回指令，但它对返回地址、标志和段状态有特定约束。完整的
中断式返回路径能从显式帧中恢复更多体系结构状态，成本也更高。两者的合理关系是：

\begin{enumerate}
\tightlist
\item 先建立一条能验证并完整恢复状态的通用返回路径；
\item 证明当前入口类型、用户地址、标志和跟踪状态满足快速路径条件；
\item 只有证明成立时才选择快速返回；
\item 任一条件不确定，都退回通用路径，而不是补猜缺失状态。
\end{enumerate}

\begin{longtable}{|L{0.20\linewidth}|L{0.29\linewidth}|L{0.39\linewidth}|}
\hline
返回情形 & 快速路径是否合适 & 理由 \\ \hline
普通 64 位调用，无信号和跟踪 & 可能合适 & 返回地址和标志保持在快速指令约束内 \\ \hline
信号改写了用户继续状态 & 需要重新证明 & 新 \code{RIP/RSP} 来自刚构造的用户现场 \\ \hline
调试器要求单步 & 通常走慢路径 & 调试标志和退出观察需要完整处理 \\ \hline
兼容 ABI 调用 & 使用对应 ABI 的路径 & 地址宽度、段状态和参数布局不同 \\ \hline
信号返回恢复扩展状态 & 通用验证路径 & 用户帧包含大量不可信恢复字段 \\ \hline
\end{longtable}

快速路径的收益来自缩短已经正确的公共路径，而不是省略语义。若优化改变了信号、跟踪或
错误结果，它就不再是同一接口的实现。

\subsection{返回前如何避免泄露内核临时状态}

系统调用 ABI 规定哪些寄存器保存、哪些可破坏，但“可破坏”不等于可以携带内核秘密。
入口代码、C 调用约定和体系结构代码会使用临时寄存器；如果返回指令恢复的集合少于内核
使用的集合，未覆盖寄存器可能暴露指针、栈内容或地址布局。

安全出口需要为每个返回寄存器给出明确来源：从用户入口帧恢复、写入规定结果、写入 ABI
允许的固定值，或在返回前清零。浮点、向量、调试和控制寄存器也要遵守 task 所有权切换
规则。不能用“用户大概不会读取”代替状态定义。

\begin{center}
\begin{tikzpicture}[
  >=Latex,
  node distance=8mm and 8mm,
  box/.style={draw=BookRule, rounded corners=2pt, align=center,
              minimum width=2.9cm, minimum height=1.0cm, fill=BookBlueTint},
  arr/.style={->, thick, draw=BookAccent}
]
\node[box] (result) {对象实现结束\\结果已确定};
\node[box, right=of result] (work) {结清 task 工作\\调度、信号、跟踪};
\node[box, right=of work] (check) {验证返回现场\\选择快/慢路径};
\node[box, right=of check] (user) {恢复用户状态\\执行下一条指令};
\draw[arr] (result) -- (work);
\draw[arr] (work) -- (check);
\draw[arr] (check) -- (user);
\draw[arr, bend right=28] (work.south) to node[below, font=\small]{新工作到达则重新检查} (work.south west);
\end{tikzpicture}
\end{center}
\diagramnote{“处理函数返回”与“处理器回到用户态”是两个时刻。中间的返回工作路径结清
task 级状态，并证明最终用户现场可安全恢复。}

\begin{keyidea}
系统调用出口不是入口的机械逆过程。入口把不可信请求转换成内核拥有的状态；出口要先结清
并发到达的 task 级工作，再把受控结果转换成一个可安全恢复的用户现场。快速返回只能建立
在完整语义已经可用、且本次状态满足额外约束的证明之上。
\end{keyidea}

\section{正确路径建立后怎样判断入口成本是否值得优化}

\subsection{先把一项只读请求的固定成本逐项记账}

读取时间、查询 CPU 编号或获得进程标识一类请求可能只需要几个字段。若每次都经历特权
切换，固定边界成本可能大于真正的数据读取成本。要讨论优化，先把基础方案的工作拆开：

\begin{longtable}{|L{0.21\linewidth}|L{0.29\linewidth}|L{0.38\linewidth}|}
\hline
基础步骤 & 即使结果很小仍需完成的工作 & 保护的语义 \\ \hline
建立入口现场 & 保存继续位置并切到内核拥有的栈 & 可恢复性与隔离 \\ \hline
入口过滤与分发 & 识别 ABI、编号和边界策略 & 稳定契约与安全策略 \\ \hline
取得共享状态 & 读取时间源或 task 字段 & 一致性 \\ \hline
返回工作检查 & 检查信号、调度和跟踪 & task 级语义 \\ \hline
恢复特权级 & 验证并恢复用户现场 & 安全交还执行权 \\ \hline
\end{longtable}

这些工作对修改全局状态的调用是必要代价。对于“频繁读取、极少修改、允许只在用户态
获得快照”的信息，重复跨越边界就值得重新审视。优化问题由此变得具体：能否把一份经过
约束的只读状态映射给用户，让用户在不进入内核的情况下得到与系统调用相同的结果？

\subsection{共享只读数据为何仍需要一致性协议}

直接映射一页时间数据并不充分。内核可能正好在用户读取一半时更新基准时间、周期换算系数
和时钟源编号；把新旧字段拼在一起会得到不存在的时间点。读者又不能取得会阻塞内核更新
的普通锁。

一种适合“单个写者协议、读者短暂读取”的安排是版本序列。写者更新前把序号改成奇数，
更新完再改成下一个偶数；读者只有在前后读到同一个偶数时才接受快照：

\begin{lstlisting}
read_shared_time:
  repeat:
    begin = shared.sequence
    if begin is odd:
      continue
    base       = shared.base_time
    cycle_base = shared.cycle_base
    multiplier = shared.multiplier
    shift      = shared.shift
    now_cycle  = read_hardware_counter()
    end = shared.sequence
  until begin == end and end is even
  return base + scale(now_cycle - cycle_base, multiplier, shift)
\end{lstlisting}

读者没有取得所有权；它靠验证“读取期间没有写者完成或正在写”来接受快照。若验证失败，
只是重读，不带着混合状态计算结果。共享页还必须只向用户暴露接口所需字段，不能顺便映射
内核内部对象。

\subsection{用户态辅助入口如何保持系统调用的后备语义}

仅有数据页还不够。不同处理器的计数器读取方式、时钟源能力和换算规则可能变化，需要一段
由内核按 ABI 提供、但在用户态执行的辅助代码。这就是从上述限制得到的 vDSO 类接口：
调用形式看起来像普通库函数；适合时在用户态读共享状态，不适合时退回真正系统调用。

\begin{center}
\begin{tikzpicture}[
  >=Latex,
  node distance=10mm and 13mm,
  box/.style={draw=BookRule, rounded corners=2pt, align=center,
              minimum width=3.25cm, minimum height=1.0cm, fill=BookBlueTint},
  arr/.style={->, thick, draw=BookAccent}
]
\node[box] (call) {用户请求时间};
\node[box, right=of call] (vdso) {辅助代码检查\\共享页与时钟源};
\node[box, above right=8mm and 15mm of vdso] (fast) {用户态计算\\版本验证成功};
\node[box, below right=8mm and 15mm of vdso] (slow) {系统调用后备\\内核完成读取};
\node[box, right=35mm of vdso] (same) {同一 ABI 结果};
\draw[arr] (call) -- (vdso);
\draw[arr] (vdso) -- node[above left, font=\small]{可安全读取} (fast);
\draw[arr] (vdso) -- node[below left, font=\small]{能力不足或重读失败} (slow);
\draw[arr] (fast) -- (same);
\draw[arr] (slow) -- (same);
\end{tikzpicture}
\end{center}
\diagramnote{用户态快速读取不是另一套时间语义。它必须与系统调用共享结果约定，并保留
无法安全完成时的内核后备路径。}

这种方法只适合不需要执行权限检查、修改对象或产生不可逆副作用的请求。若把“打开文件”
或“改变系统时间”也做成用户态共享写入，内核就失去了验证与提交边界。

\subsection{为什么减少调用次数和缩短单次入口是两种优化}

设一次边界往返固定成本为 \(C_e\)，每项请求的实际处理成本为 \(C_o\)，连续提出 \(n\)
项请求的基础成本近似为：

\[
T_{\mathrm{separate}} = n(C_e + C_o)
\]

若协议允许一次提交 \(n\) 项独立描述符，边界成本可以摊薄：

\[
T_{\mathrm{batch}} = C_e + nC_o + C_q
\]

其中 \(C_q\) 是队列解析、描述符验证和完成记录成本。批处理优化的是入口次数；vDSO 类
方案优化的是某一类只读请求的整个入口。二者不能互相替代。

批处理也不是把多个原调用简单拼接。接口必须说明：

\begin{itemize}
\tightlist
\item 描述符由谁拥有，提交后用户能否改写；
\item 一项失败是否阻止后续项，还是每项独立产生结果；
\item 完成顺序是否等于提交顺序；
\item 部分完成时怎样报告已提交前缀；
\item task 退出、对象关闭或信号到达时怎样取消；
\item 完成记录何时对用户可见，是否可能覆盖未消费结果。
\end{itemize}

如果这些问题没有定义，减少入口次数只会把同步错误从单个调用移到共享队列。

\subsection{异步提交为什么需要把调用栈中的状态搬到独立对象}

同步调用期间，参数快照、对象引用和进度可以位于当前 task 的内核栈上，因为完成发生在
该调用返回之前。异步提交允许 task 先返回，真正操作稍后由另一个执行上下文完成。原内核
栈会被复用，局部变量不再存在；于是操作状态必须成为独立、可引用的请求对象。

\begin{longtable}{|L{0.21\linewidth}|L{0.27\linewidth}|L{0.40\linewidth}|}
\hline
请求对象字段组 & 保存内容 & 存在理由 \\ \hline
身份与关联 & 请求编号、提交者、完成队列 & 让完成结果找到原请求 \\ \hline
规范化参数 & 长度、偏移、操作码、固定缓冲描述 & 用户描述符提交后可能改变 \\ \hline
对象引用 & 文件、套接字、内存注册对象 & task 返回或关闭槽位后仍要保证生命期 \\ \hline
进度与结果 & 已提交字节、错误、完成状态 & 支持部分完成和唯一终态 \\ \hline
取消状态 & 取消请求、不可取消点、竞争结果 & 定义取消与完成同时发生时的唯一答案 \\ \hline
凭据与策略上下文 & 规定时刻捕获的身份或引用 & 后台执行不能随意使用工作线程身份 \\ \hline
\end{longtable}

这一步揭示一个重要差异：异步接口不是“系统调用返回得更快”，而是改变了请求状态的所有权
和生命期。只把用户指针存入队列、稍后直接访问，会重新引入地址变化、页面回收和 TOCTOU
问题。实现要么在提交时复制小参数，要么注册并固定受约束的缓冲区，并规定解除注册与在途
请求的同步关系。

\begin{center}
\begin{tikzpicture}[
  >=Latex,
  node distance=10mm and 12mm,
  box/.style={draw=BookRule, rounded corners=2pt, align=center,
              minimum width=3.15cm, minimum height=1.0cm, fill=BookBlueTint},
  arr/.style={->, thick, draw=BookAccent}
]
\node[box] (desc) {用户描述符\\仍由用户控制};
\node[box, right=of desc] (snap) {提交边界\\复制或固定};
\node[box, right=of snap] (req) {内核请求对象\\引用、进度、凭据};
\node[box, below=14mm of req] (worker) {后台执行上下文\\推进唯一状态机};
\node[box, left=of worker] (cq) {完成记录\\结果与请求编号};
\draw[arr] (desc) -- (snap);
\draw[arr] (snap) -- (req);
\draw[arr] (req) -- (worker);
\draw[arr] (worker) -- (cq);
\draw[arr] (cq) -- node[left, font=\small]{用户消费后释放} (snap);
\end{tikzpicture}
\end{center}
\diagramnote{同步调用依赖当前内核栈保存进度；异步提交必须把仍需存活的状态迁移到独立请求
对象。描述符被接受、操作提交和完成可见是三个不同时间点。}

\subsection{接受、提交和完成为什么必须分开}

异步接口返回“接受成功”只说明描述符已通过初步验证并进入内核拥有的队列。设备或对象可能
尚未看到请求，更没有产生结果。至少要区分：

\begin{longtable}{|L{0.18\linewidth}|L{0.31\linewidth}|L{0.39\linewidth}|}
\hline
时刻 & 已经成立的事实 & 尚不能推断的事实 \\ \hline
接受 & 内核已取得请求状态的所有权 & 操作已到达对象或设备 \\ \hline
提交 & 对象或下层队列已经接收工作 & 数据已复制、落盘或被对端接收 \\ \hline
完成 & 接口定义的结果已经确定并发布 & 若接口只保证缓存完成，则不代表持久化 \\ \hline
用户消费 & 用户已读到完成记录 & 对象的外部副作用已经被撤销 \\ \hline
\end{longtable}

取消请求也只能相对于这些状态定义。尚未提交时通常可从队列删除；已经提交后可能只能标记
“不再需要结果”，不能保证设备停止；完成与取消并发时，状态机必须以原子转换决定究竟发布
成功、取消还是已完成。不能同时产生两个终态。

\subsection{兼容 ABI 为什么不能只截断寄存器}

同一内核可能服务不同地址宽度或历史布局的用户程序。32 位调用者传入的指针、\code{long}
和结构体对齐方式与 64 位内核不同。若只把寄存器高位清零然后调用同一实现，嵌套结构体仍
会按错误偏移解释。

\begin{longtable}{|L{0.20\linewidth}|L{0.28\linewidth}|L{0.40\linewidth}|}
\hline
差异来源 & 直接复用的风险 & 兼容层的职责 \\ \hline
指针宽度 & 高位扩展或范围判断错误 & 明确转成内核地址表示并验证用户范围 \\ \hline
\code{long}/时间字段宽度 & 溢出、符号扩展和时间范围变化 & 按旧布局取值，再转为统一内核类型 \\ \hline
结构体填充与对齐 & 从错误字节读取字段 & 使用对应 ABI 布局逐字段转换 \\ \hline
系统调用编号表 & 同一编号可能表示不同操作 & 先识别 ABI，再进入相应分发表 \\ \hline
返回值宽度 & 部分结果被静默截断 & 按接口规定检测溢出并返回兼容错误 \\ \hline
\end{longtable}

合理的共享边界位于“已经转换成内核统一表示”之后：各 ABI 包装层负责读取旧格式、检查
范围并形成规范化参数；对象实现只处理统一类型。这样既避免复制业务逻辑，也不会让核心
实现到处猜测调用者位宽。

\subsection{安全缓解措施为什么也属于边界成本}

现代处理器可能在架构检查完成前推测执行后续指令。系统调用入口跨越信任边界，返回路径
又可能改变预测上下文；某些平台需要序列化、预测器状态管理或间接分支约束。这些措施增加
固定入口成本，但它们处理的是微体系结构可观察性，不能从 C 级权限检查中推导出来。

书写和优化时应保持两个层次：

\begin{itemize}
\tightlist
\item 架构语义说明寄存器、页表、特权级和提交状态必须满足什么条件；
\item 平台缓解说明为了让推测执行不跨越该条件，还需执行哪些附加步骤。
\end{itemize}

缓解策略会随处理器能力和系统配置变化，因此不应把某一组指令写成系统调用概念本身。
无论启用何种缓解，参数快照、对象引用、授权与返回验证仍是基础正确性；关闭缓解也不能
省略它们。

\begin{keyidea}
边界优化必须从一条已正确、可计量的基础路径出发。只读共享快照减少某类入口，批处理摊薄
多项请求的固定成本，异步提交改变请求状态的所有权与生命期，兼容层把不同用户布局转换成
统一内核表示。每项优化都要保留原接口的权限、提交、失败和返回语义。
\end{keyidea}

\section{把各阶段合成一项完整而可核对的请求}

\subsection{先跟踪一次会阻塞的 pipe read}

现在才把前面逐项得到的结构连成完整路径。设 task T 调用
\code{read(fd, user\_buf, 4096)}，管道开始为空，随后写者写入 600 字节，同时 T 收到一个
待处理信号。为了避免“调用成功”和“信号到达”混成一句话，逐时刻记录所有权与提交状态。

\begin{longtable}{|L{0.10\linewidth}|L{0.25\linewidth}|L{0.25\linewidth}|L{0.28\linewidth}|}
\hline
时刻 & T 或 CPU 状态 & 关键对象状态 & 已经成立的接口事实 \\ \hline
\(t_0\) & T 在用户态 & fd 槽引用管道读端 & 尚未提出本次请求 \\ \hline
\(t_1\) & 入口帧已建立 & 参数仍只是整数与用户地址 & 可恢复，但尚未取得对象 \\ \hline
\(t_2\) & 处理函数运行 & 已取得读端引用，长度已验证 & fd 重用不再改变本次对象 \\ \hline
\(t_3\) & T 在等待队列并睡眠 & 管道为空，未提交字节 & 调用仍在内核中进行 \\ \hline
\(t_4\) & 写者唤醒 T & 管道有 600 字节，T runnable & 只获得运行资格，尚未完成读取 \\ \hline
\(t_5\) & T 恢复并持有管道锁 & 600 字节可消费，信号挂起 & 条件成立，信号尚未决定结果 \\ \hline
\(t_6\) & T 复制并推进读位置 & 600 字节已交付给用户缓冲区 & 完成前缀为 600，不能从头重启 \\ \hline
\(t_7\) & 返回工作选择信号 & 信号帧构造在用户栈 & 主返回值保持 600 \\ \hline
\(t_8\) & 处理器回到用户态 & \code{RIP} 指向信号处理函数 & 处理函数之后可读到 600 的调用结果 \\ \hline
\end{longtable}

在 \(t_2\) 之后关闭同一 fd 槽位，不会改变 T 已取得的对象引用。在 \(t_4\) 唤醒后，另一
读者可能先消费数据，所以 T 到 \(t_5\) 仍要复查。到 \(t_6\) 才首次出现接口可见提交；
从这一刻起，信号不能把主结果改成“完全未发生”的中断错误。

\subsection{600 字节究竟在哪一个时刻算作已读取}

“从管道移除”和“写入用户缓冲区”必须安排成不会无故丢失数据的协议。若先推进管道读位置，
随后用户页故障且复制无法完成，数据已经从管道消失，调用者却没有得到它。若先复制又允许
其他读者同时消费同一区间，则会重复交付。

实现需要在对象同步边界内为当前读者保留区间，按可处理的片段受控复制，并只对成功复制的
字节推进提交位置。具体锁粒度可以优化，但必须维持：

\[
0 \leq \text{committed} \leq \text{reserved} \leq
\min(\text{requested},\text{available})
\]

用户复制在第 300 字节失败时，接口可返回 300，而对象只消费对应的 300 字节；尚未成功
复制的保留区间要释放给后续读取。底层复制函数报告的“剩余 300”要先换算成“成功 300”，
再由系统调用层根据是否存在成功前缀选择主结果。

\subsection{再跟踪一次改变全局时钟的请求}

回到本章开头的问题。管理员程序请求把系统实时时钟改到新值。与读取管道不同，这项操作
没有可自然返回的字节前缀，而且会影响全局观察者。完整路径应当说明：

\begin{enumerate}
\item 用户库把操作编号、时钟类型和用户结构地址放入 ABI 规定位置。
\item 入口代码保存用户继续状态，切换到 task 的内核栈，建立统一入口帧。
\item 分发层先完成入口过滤，再由包装层从用户结构取得时间值并转换成统一内核表示。
\item 转换过程检查字段范围，例如纳秒必须位于接口允许区间；用户随后改写原结构不再影响
      已取得的内核快照。
\item 授权层使用当前 task 的凭据，检查相对于正确 user namespace 的时间管理能力。
\item 时间子系统在自己的写侧同步边界内重新读取相关状态，计算并提交新基准；授权成功并不
      预留未来提交权。
\item 提交成功后更新保证用户只读快照一致性的版本序列，并通知需要观察时间变化的子系统。
\item 对象实现写入零结果；返回工作处理调度、信号、审计和跟踪，最后验证用户现场。
\end{enumerate}

如果权限检查失败，第 6 步不能发生；如果用户结构无效，甚至不应进入授权；如果提交前参数
范围错误，时钟保持原值。审计记录要区分“请求过什么值”“以哪份凭据判断”“是否真正提交”
与“最终向用户返回什么结果”。

\begin{center}
\begin{tikzpicture}[
  >=Latex,
  node distance=7mm,
  box/.style={draw=BookRule, rounded corners=2pt, align=left,
              minimum width=10.3cm, minimum height=0.76cm, fill=BookBlueTint},
  arr/.style={->, thick, draw=BookAccent}
]
\node[box] (a) {1.\ 入口：保存用户继续状态，只得到编号与原始参数};
\node[box, below=of a] (b) {2.\ 快照：复制时间结构，检查布局、范围和 ABI};
\node[box, below=of b] (c) {3.\ 授权：用当前凭据判断对目标时间域的操作权};
\node[box, below=of c] (d) {4.\ 提交：在时间子系统同步边界内重读状态并更新基准};
\node[box, below=of d] (e) {5.\ 发布：更新共享读取序列，形成审计结果};
\node[box, below=of e] (f) {6.\ 返回：结清 task 工作，验证并恢复用户现场};
\draw[arr] (a) -- (b);
\draw[arr] (b) -- (c);
\draw[arr] (c) -- (d);
\draw[arr] (d) -- (e);
\draw[arr] (e) -- (f);
\end{tikzpicture}
\end{center}
\diagramnote{这条链路直到本章末才完整出现，因为编号、入口帧、用户快照、凭据、提交与返回
边界已经分别由具体问题推导出来。}

\subsection{失败注入怎样验证边界真的存在}

正常执行只能证明某条路径可用，不能证明状态在失败时仍闭合。可以在不同边界人为制造失败，
逐项检查“谁拥有状态、是否已经提交、应返回什么”：

\begin{longtable}{|L{0.20\linewidth}|L{0.30\linewidth}|L{0.38\linewidth}|}
\hline
注入位置 & 预期保持的状态 & 应观察的结果 \\ \hline
入口编号不存在 & 不取得对象、不复制用户缓冲区 & 返回未实现类错误，用户继续状态完整 \\ \hline
用户缓冲区第二页缺失 & 只提交成功复制的前缀 & 返回前缀数量；若前缀为零则返回地址错误 \\ \hline
fd 在取得引用前被关闭 & 不得使用已释放对象 & 取得槽位同步决定旧对象或无效 fd \\ \hline
fd 在取得引用后被关闭 & 当前操作仍持有对象生命期 & 槽位可重用，本次请求继续引用原对象 \\ \hline
权限在提交前变化 & 按接口规定时刻重新判断或用快照 & 不把曾经检查过等同于永久授权 \\ \hline
等待时收到信号 & 保留已提交前缀和对象一致性 & 无前缀时中断或重启，有前缀时返回数量 \\ \hline
信号帧用户栈不可写 & 不恢复到半构造现场 & 进入规定的失败或终止路径 \\ \hline
快速返回条件被破坏 & 完整状态仍可恢复 & 退回通用返回路径，结果语义不变 \\ \hline
\end{longtable}

每个测试都应同时观察接口结果和内核对象状态。只断言“返回了 \code{-EFAULT}”不够；还要
确认没有泄漏对象引用、没有重复消费数据、没有发布半初始化 fd，也没有把错误的用户现场
交给处理器。

\subsection{用五个问题审查任意新系统调用}

面对一个新的系统调用，不必先背诵完整内核链路。按状态形成顺序提出五个问题即可定位主要
结构：

\begin{enumerate}
\item \textbf{请求如何成为稳定输入？} 哪些标量可直接取值，哪些用户结构必须复制，嵌套
      指针和尺寸乘法怎样验证？
\item \textbf{操作对象怎样获得稳定身份？} 是从 fd、路径还是其他句柄取得；引用何时增加，
      槽位关闭或重用是否会影响本次操作？
\item \textbf{谁在什么时刻授权哪项操作？} 使用当前、捕获还是对象凭据；检查依赖哪些对象
      状态，提交前是否必须重读？
\item \textbf{哪一步首次产生不可忽略的外部效果？} 部分完成怎样计数；失败和信号能否重启；
      发布、完成和持久化是否为不同边界？
\item \textbf{怎样形成安全的用户继续状态？} 返回前有哪些 task 工作；输出是否完整初始化；
      用户地址、标志和扩展状态由谁验证？
\end{enumerate}

这五问不是固定写作模板，而是一组状态完整性检查。对于 \code{getpid}，第二问几乎为空；
对于异步 I/O，第二和第四问会展开成独立请求对象与状态机；对于信号返回，第五问本身就是
主要实现。问题的篇幅由机制决定，而不是强行平均。

\subsection{本章最终得到的状态分层}

从“用户不能直接改全局时钟”出发，本章没有立即假设一个完整操作系统入口。现在可以把已经
逐步得到的状态按所有者归纳：

\begin{longtable}{|L{0.18\linewidth}|L{0.31\linewidth}|L{0.39\linewidth}|}
\hline
所有者 & 代表状态 & 负责回答的问题 \\ \hline
处理器与每 CPU 入口区 & 入口地址、最早栈信息、CPU 局部状态 & 怎样安全离开用户执行环境 \\ \hline
task 与内核栈 & 入口帧、调用链、切换帧、返回工作标志 & 本次控制流在哪里继续 \\ \hline
系统调用边界层 & ABI、编号、包装与过滤状态 & 用户究竟请求了什么 \\ \hline
用户快照或固定映射 & 规范化参数、输出暂存、页引用 & 用户内存变化是否还能影响操作 \\ \hline
对象与引用 & 打开对象、路径对象、请求对象 & 操作作用于谁，生命期到何时 \\ \hline
凭据与安全策略 & 身份、能力集合、命名空间、LSM 上下文 & 谁能对该对象做哪项操作 \\ \hline
等待与调度结构 & 等待队列、runqueue、已提交前缀 & 条件未满足时怎样停下并继续 \\ \hline
返回协议 & 信号帧、重启状态、快慢路径条件 & 怎样把结果安全交还用户 \\ \hline
\end{longtable}

这些结构并不都在每次调用中出现。它们分别处理不同的不稳定性：控制流会被打断，用户内存
会变化，fd 槽会重用，凭据与对象状态会并发改变，等待会跨越调度，返回现场会被信号重写。
结构存在的理由就是封闭相应变化；若某项调用不接触那类变化，就不应为了形式完整而制造
无用状态。

\begin{keyidea}
一次系统调用可以概括为四次所有权变化：处理器把用户继续状态交给 task 的入口帧；边界层
把用户控制的参数变成内核拥有的快照；对象层取得稳定引用并在明确时刻提交效果；返回层把
内核结果变成经过验证的用户继续状态。阻塞、信号、跟踪和优化都不能破坏这四次转换。
\end{keyidea}
